%~Mouliné par MaN_auto v.0.37.1 (a7a71208) 2025-10-28 10:29:44
\documentclass[AHL,Unicode,longabstracts]{cedram}








\newcommand{\pG}{\partial{\Gamma}}
\newcommand{\dirac}[1]{{\mathrm{Dirac}}{(#1)}}



\newcommand{\N}{\mathbb{N}}
\newcommand{\Z}{\mathbb{Z}}
\newcommand{\bbP}{\mathbb{P}}
\newcommand{\R}{\mathbb{R}}
\newcommand\bbN{\mathbb{N}}
\newcommand{\rme}{\mathrm{e}}
\newcommand{\rmd}{\mathrm{d}}

\newcommand{\fkd}{\mathfrak{d}}
\newcommand{\fkB}{\mathfrak{B}}
\newcommand{\bfE}{\mathbf{E}}
\newcommand{\bfP}{\mathbf{P}}

\newcommand{\clM}{\mathcal{M}}
\newcommand{\clP}{\mathcal{P}}
\newcommand{\clV}{\mathcal{V}}
\newcommand{\clZ}{\mathcal{Z}}
\newcommand{\clW}{\mathcal{W}}


\DeclareMathOperator{\BRW}{BRW}
\DeclareMathOperator{\supp}{supp}
\DeclareMathOperator{\con}{con}
\DeclareMathOperator{\Cay}{Cay}
\DeclareMathOperator{\sd}{sd}


%%%%%%%%%%%%%%%%%%%%%%%%%%%%%%%%%%%%%%%%%%%%%%%%%%%%%%%%%%%%%%%%%%%%%%%%%%%%%%%%%%%%
%Better widebar

\makeatletter
\let\save@mathaccent\mathaccent
\newcommand*\if@single[3]{%
  \setbox0\hbox{${\mathaccent"0362{#1}}^H$}%
  \setbox2\hbox{${\mathaccent"0362{\kern0pt#1}}^H$}%
  \ifdim\ht0=\ht2 #3\else #2\fi
  }
%The bar will be moved to the right by a half of \macc@kerna, which is computed by amsmath:
\newcommand*\rel@kern[1]{\kern#1\dimexpr\macc@kerna}
%If there's a superscript following the bar, then no negative kern may follow the bar;
%an additional {} makes sure that the superscript is high enough in this case:
\newcommand*\widebar[1]{\@ifnextchar^{{\wide@bar{#1}{0}}}{\wide@bar{#1}{1}}}
%Use a separate algorithm for single symbols:
\newcommand*\wide@bar[2]{\if@single{#1}{\wide@bar@{#1}{#2}{1}}{\wide@bar@{#1}{#2}{2}}}
\newcommand*\wide@bar@[3]{%
  \begingroup
  \def\mathaccent##1##2{%
%Enable nesting of accents:
    \let\mathaccent\save@mathaccent
%If there's more than a single symbol, use the first character instead (see below):
    \if#32 \let\macc@nucleus\first@char \fi
%Determine the italic correction:
    \setbox\z@\hbox{$\macc@style{\macc@nucleus}_{}$}%
    \setbox\tw@\hbox{$\macc@style{\macc@nucleus}{}_{}$}%
    \dimen@\wd\tw@
    \advance\dimen@-\wd\z@
%Now \dimen@ is the italic correction of the symbol.
    \divide\dimen@ 3
    \@tempdima\wd\tw@
    \advance\@tempdima-\scriptspace
%Now \@tempdima is the width of the symbol.
    \divide\@tempdima 10
    \advance\dimen@-\@tempdima
%Now \dimen@ = (italic correction / 3) - (Breite / 10)
    \ifdim\dimen@>\z@ \dimen@0pt\fi
%The bar will be shortened in the case \dimen@<0 !
    \rel@kern{0.6}\kern-\dimen@
    \if#31
      \overline{\rel@kern{-0.6}\kern\dimen@\macc@nucleus\rel@kern{0.4}\kern\dimen@}%
      \advance\dimen@0.4\dimexpr\macc@kerna
%Place the combined final kern (-\dimen@) if it is >0 or if a superscript follows:
      \let\final@kern#2%
      \ifdim\dimen@<\z@ \let\final@kern1\fi
      \if\final@kern1 \kern-\dimen@\fi
    \else
      \overline{\rel@kern{-0.6}\kern\dimen@#1}%
    \fi
  }%
  \macc@depth\@ne
  \let\math@bgroup\@empty \let\math@egroup\macc@set@skewchar
  \mathsurround\z@ \frozen@everymath{\mathgroup\macc@group\relax}%
  \macc@set@skewchar\relax
  \let\mathaccentV\macc@nested@a
%The following initialises \macc@kerna and calls \mathaccent:
  \if#31
    \macc@nested@a\relax111{#1}%
  \else
%If the argument consists of more than one symbol, and if the first token is
%a letter, use that letter for the computations:
    \def\gobble@till@marker##1\endmarker{}%
    \futurelet\first@char\gobble@till@marker#1\endmarker
    \ifcat\noexpand\first@char A\else
      \def\first@char{}%
    \fi
    \macc@nested@a\relax111{\first@char}%
  \fi
  \endgroup
}
\makeatother

\let\oldbar\bar
\renewcommand*{\bar}[1]{{\mathchoice{\widebar{#1}}{\widebar{#1}}{\widebar{#1}}{\oldbar{#1}}}}
%%%%%%%%%%%%%%%
%~ This fixes spacing issues around \left( and \right) and other similar symbols
\let\originalleft\left 
\let\originalright\right
\renewcommand{\left}{\mathopen{}\mathclose\bgroup\originalleft}
\renewcommand{\right}{\aftergroup\egroup\originalright}
%%%%%%%%
\let\oldin\in
\renewcommand{\in}{\mathchoice{\oldin}{\oldin}{\,\oldin\,}{\,\oldin\,}}

\let\oldge\ge
\renewcommand{\ge}{\mathchoice{\oldge}{\oldge}{\,\oldge\,}{\,\oldge\,}}

\let\oldle\le
\renewcommand{\le}{\mathchoice{\oldle}{\oldle}{\,\oldle\,}{\,\oldle\,}}

\renewcommand*{\to}{\mathchoice{\longrightarrow}{\rightarrow}{\rightarrow}{\rightarrow}}

\let\oldmapsto\mapsto
\renewcommand*{\mapsto}{\mathchoice{\longmapsto}{\oldmapsto}{\oldmapsto}{\oldmapsto}}
%%%%%%%%%%%%%%%%%%%%%%%%%%%%%%%%%%%%%%%%%%%%%%%%%%%%%%%%%%%%%%%%%%%%%%%%%%%%%%%%%%%%

\graphicspath{{./figures/}}

\newcommand*{\mk}{\mkern -1mu}
\newcommand*{\Mk}{\mkern -2mu}
\newcommand*{\mK}{\mkern 1mu}
\newcommand*{\MK}{\mkern 2mu}

\hypersetup{urlcolor=purple, linkcolor=blue, citecolor=red}

\newcommand*{\relabel}{\renewcommand{\labelenumi}{(\theenumi)}}
\newcommand*{\romanenumi}{\renewcommand*{\theenumi}{\roman{enumi}}\relabel}
\newcommand*{\Romanenumi}{\renewcommand*{\theenumi}{\Roman{enumi}}\relabel}
\newcommand*{\alphenumi}{\renewcommand*{\theenumi}{\alph{enumi}}\relabel}
\newcommand*{\Alphenumi}{\renewcommand*{\theenumi}{\Alph{enumi}}\relabel}
\let\oldhat\hat
\renewcommand*{\hat}[1]{\mathchoice{\widehat{#1}}{\widehat{#1}}{\oldhat{#1}}{\oldhat{#1}}}
\let\oldexists\exists
\renewcommand*{\exists}{\mathrel{\oldexists}}
\let\oldforall\forall
\renewcommand*{\forall}{\mathrel{\oldforall}}

\newcommand{\llbangle}{\mathopen{\langle\mkern -3.7mu\langle}}
\newcommand{\rrbangle}{\mathclose{\rangle\mkern -3.7mu\rangle}}
%%%%%%%%%%%%%%%%%%%%%%%%%%%%%%%%%%%%%%%%%%%%%%%%%%%%%%%%%%%%%%%%%%%%%%%%%%%%%%%%%%%%

\title[The growth of the Green function and Poincaré series]{The growth of the Green function for random walks and Poincaré series}
\alttitle{Croissance de la fonction de Green pour les marches aléatoires et séries de Poincaré}

\subjclass{20F65, 20P05, 60J80}
\keywords{Branching random walks, relatively hyperbolic group, growth rate, Green function, Poincaré series, Growth tightness, Patterson--Sullivan measure, parabolic gap}


\author[\initial{M.} \lastname{Dussaule}]{\firstname{Matthieu} \lastname{Dussaule}}
\address{Laboratoire Cogitamus}
\email{matthieu.dussaule@hotmail.fr}





\author[\initial{L.} \lastname{Wang}]{\firstname{Longmin} \lastname{Wang}}
\address{School of Statistics\\
and Data Science, \\
LPMC, Nankai University,\\
94 Weijin Road,\\
Nankai District, Tianjin (China)}
\email{wanglm@nankai.edu.cn}

\author[\initial{W.} \lastname{Yang}]{\firstname{Wenyuan} \lastname{Yang}}
\address{Beijing International Center\\
for Mathematical Research (BICMR),\\
Pekin University,\\
No. 5 Yiheyuan Road,\\
Haidian District, Beijing (China)}
\email{yabziz@gmail.com}

\thanks{L.~W. is supported by National Natural Science Foundation of China (12171252). W.~Y. is supported by National Key R \& D Program of China (SQ2020YFA070059) and National Natural Science Foundation of China (no.~12131009 and no~12326601).}
\CDRGrant[National Natural Science Foundation of China]{12171252}
\CDRGrant[National Key R \& D Program of China]{SQ2020YFA070059}
\CDRGrant[National Natural Science Foundation of China]{12131009}
\CDRGrant[National Natural Science Foundation of China]{12326601}

\begin{abstract}
Given a probability measure $\mu$ on a finitely generated group $\Gamma$, the Green function $G(x,y|r)$ encodes many properties of the random walk associated with $\mu$. Endowing $\Gamma$ with a word distance, we denote by $H_r(n)$ the sum of the Green function $G(e,x|r)$ along the sphere of radius $n$. This quantity appears naturally when studying asymptotic properties of branching random walks driven by $\mu$ on $\Gamma$. Our main result exhibits a relatively hyperbolic group with convergent Poincaré series generated by $H_r(n)$.
\end{abstract}

\begin{altabstract}
Étant donné une mesure de probabilité $\mu$ sur un groupe de type fini $\Gamma$, la fonction de Green $G(x,y|r)$ encore de nombreuses propriétés de la marche aléatoire associée à $\mu$. Lorsque $\Gamma$ est muni d'une distance des mots, on note $H_r(n)$ la somme de la fonction de Green le long de la sphère de rayon $n$. Cette quantité apparaît de manière naturelle dans l'étude des propriétés asymptotiques des marches aléatoires branchantes dirigées par $\mu$ sur $\Gamma$. Le résultat principal exhibe un groupe relativement hyperbolique dont la série de Poincaré engendrée par $H_r(n)$ est convergente.
\end{altabstract}




\datereceived{2023-07-21}
\daterevised{2025-09-10}
\dateaccepted{2025-09-17}

\editors{C. Sorger and V. Guirardel}
%%%%%%%%%%%%%%%%%%%%%%%%%%%%%%%%%%%%%%%%%%%%%%%%%%%%%%%%%%%%%%%%%%%%%%%%%%%%%%%%%%%%
\begin{document}

\maketitle

\section{Introduction}


Let $\Gamma$ be a finitely generated group endowed with a finite generating set. Denote by $|\!\cdot\!|$ the induced word distance and by $S_n = \left\{ x \in \Gamma \colon |x| =n \right\}$ the sphere of radius~$n$ centered at the neutral element $e$ of $\Gamma$. Set $v=\limsup_{n\to\infty} \frac{1}{n}\log \sharp S_n$. The number~$v$ is called the volume growth of $\Gamma$ with respect to the chosen generating set. The group $\Gamma$ is said to have exponential growth if $v>0$.


Consider a probability measure $\mu$ on $\Gamma$ and assume that $\mu$ is admissible, i.e. its support generates $\Gamma$ as a semigroup. Let $p_n(x, y)$ be the $n$-step transition probability for the random walk $(X_k)$ with step distribution $\mu$. Define for $0 < r \leq R$ the Green function
\[
G(x, y | r) = \sum_{n = 0}^{\infty} r^n p_n(x, y),
\]
where $R$ is the radius of convergence of the series, which is also the inverse of the spectral radius of the random walk. When $r = 1$ we write $G(x, y) = G(x, y | 1)$ for simplicity. Note that we have $G(x, y | r) = G(e, e | r) F(x, y | r)$, where $F(x, y | r)$ is the first return Green function~\cite[Lemma~1.13]{Woessbook}. Define the \emph{growth function} associated with the Green function $H_r(n)$ as
\[
H_r(n) = \sum_{x \in S_n} G(e, x | r)
\]
and the \emph{growth rate of the Green function} $\omega_\Gamma(r)$ as
\[
\omega_{\Gamma}(r) = \limsup_{n \to \infty} \frac{1}{n} \log H_r(n).
\]

We will always consider admissible random walks that are transient at the spectral radius, i.e. $G(x,y|r)$ is finite for every $x,y\in \Gamma$ and for every $r\leq R$. By Varopoulos theorem~\cite[Theorem~7.8]{Woessbook}, only groups with at most quadratic growth can carry a random walk which is not transient at the spectral radius. In other words, we will always assume that $\Gamma$ is not virtually $\Z^d$, $d\leq 2$. In particular, this ensures that $H_r(n)$ and $\omega_\Gamma(r)$ are well defined for every $r\leq R$.

\subsection*{Notations}
In all the paper, given two functions $f$ and $g$, we write $f\simeq g$ if the difference between $f$ and $g$ is bounded from above and below, that is $|f-g|\leq C$ for some constant $C$. We write $f\simeq_L g$ if the constant $C$ depends on a parameter $L$. Assuming further that $f$ and $g$ are positive, we write $f\asymp g$ if the ratio of $f$ and $g$ is bounded\pagebreak{} from above and below, that is $\frac{1}{C}f\leq g\leq Cf$ for some positive constant $C$. Similarly, if $C$ depends on $L$, we write $f\asymp_L g$. Also, if $f\leq Cg$, we write $f\prec g$ and $f\prec_L g$ if $C$ depends on $L$. If the dependency is not clear from the context, we will avoid these notations.


This paper mostly deals with the asymptotics of $H_r(n)$ as $n$ goes to infinity. A~fruitful line of research in the theory of random walks is to compute asymptotics in space of the Green function, that is asymptotics of $G(e,x|r)$ as $x$ goes to infinity. This is referred as renewal theory and goes back to Blackwell's renewal theorem for drifted random walks on $\R$, see~\cite{Blackwell} and earlier references therein. The terminology renewal comes from an interpretation of the Green function $G(e,x)$ as the probability that a renewal event takes place at $x$ for a suitable process, see~\cite[Chapter~II.9]{Spitzer}.

In fact, the terminology renewal theory is used in a much broader setting and we refer to~\cite[Chapter~XI]{Feller} for a more complete exposition within the scope of probability theory. Let us also mention that Lalley~\cite{Lalleyrenewal} generalized classical renewal theory, with a new approach to deal with asymptotics of certain counting functions arising in geometric group theory. This led to significant research in dynamical systems, where renewal theory is now a common thread.

Finally, note that Ledrappier interpreted the computation of the limit of $H_1(n)$ on the free group as a renewal theorem for the distance~\cite{Ledrappier01}, see also~\cite{Ledrappier93} for related results concerning the Brownian motion on the universal cover of compact negatively curved manifolds. However, besides such specific examples, the behaviors of $H_r(n)$ and $\omega_\Gamma(r)$ have not been investigated much in literature, so let us first explain in which context these quantities occur.



Consider a probability measure $\nu$ on $\Z_{\ge 0}=\{0,1,2,\dots\}$. The branching random walk driven by $\nu$ and $\mu$ on $\Gamma$, denoted by $\BRW(\Gamma,\nu,\mu)$ is described as follows. One starts with a single particle at $e$. For every $n$, every alive particle at time $n$ dies after giving birth to an independent number of children, according to the distribution of~$\nu$, each of which independently moves on $\Gamma$ according to the distribution of $\mu$. The~measure $\nu$ is called the offspring distribution.

In recent years, there has been a large body of work dedicated to understanding the asymptotic behavior of $\BRW(\Gamma,\nu,\mu)$ in terms of geometric features of $\Gamma$. For such study, one usually condition on non-extinction of the system, which boils down to considering an offspring measure distributed on $\N=\{1,2,\dots\}$, see~\cite[Chapter~1]{AN04}. It is thus natural to assume that $\bfE[\nu]>1$, otherwise, conditioned on non-extinction, $\nu$ is the Dirac measure at 1 and the branching random walk is nothing but the usual random walk whose step distribution is given by $\mu$.

One of the cornerstone result is that letting $\bfE[\nu]=r$, then $1<r\leq R$ if and only if the branching random walk is transient, i.e. almost surely it does not visit every vertex infinitely many times, see~\cite{BP94,GM06}. Furthermore, if $1<r\leq R$, then the branching random walk has exponential volume growth by~\cite{BM12}.

Precisely, letting $M_n$ be the cardinality of the number of elements of $S_n$ that are ever visited by the branching random walk, we have that almost surely,
\[
\limsup_{n\to\infty} \frac{1}{n}\log M_n>0.
\]


In some cases, such as hyperbolic groups~\cite[Theorem~1.1]{SWX} and relatively hyperbolic groups~\cite[Theorem~1.1]{DWY}, it was shown that this growth rate coincides almost surely with the growth rate of the Green function at $1<r\le R$. In~other words, almost surely,
\[
\limsup_{n\to\infty} \frac{1}{n}\log M_n=\omega_\Gamma(r).
\]
Moreover, for hyperbolic groups, we have $H_r(n)\asymp \rme^{n\omega_\Gamma(r)}$, i.e. the Green function has purely exponential growth, see~\cite[Theorem~3.1]{SWX}.



In fact, our original motivation for this paper was to answer some questions raised in~\cite{DWY} about branching random walks on relatively hyperbolic groups. Recall that a finitely generated group is called relatively hyperbolic if it acts via a geometrically finite action on a proper geodesic hyperbolic metric space; see Section~\ref{subsectiondefRHG} for more details. We introduce the following Poincaré series associated with $\mu$ on $\Gamma$. For $r\leq R$ and $s\in \R$, we set
\begin{align}\label{PoincareSeriesEQ}
\Theta_\Gamma(r,s)=\sum_{x\in \Gamma}G(e,x|r)\rme^{-\sd(e,x)}=\sum_{n\ge 0}H_r(n)\rme^{-sn}.
\end{align}
The growth rate $\omega_\Gamma(r)$ is thus the critical exponent of this Poincaré series, i.e. for $s<\omega_\Gamma(r)$, $\Theta_\Gamma(r,s)$ diverges and for $s>\omega_\Gamma(r)$, $\Theta_\Gamma(r,s)$ converges.

Several criteria were found in~\cite{DWY} to ensure that this Poincaré series diverges at $s=\omega_\Gamma(r)$. One of the missing pieces was whether there exists an example with convergent Poincaré series. Our main result exhibits such an example by actually constructing a relatively hyperbolic group for which divergence of the Poincaré series depends on $r$. Precisely, we prove the following.

\begin{theo}[{Theorem~\ref{transitionphasePoincare}}]\label{mainthm}
There exists a relatively hyperbolic group $\Gamma$, endowed with a finitely supported symmetric admissible probability measure $\mu$, and there exist $1<r_*<r_\sharp<R$ such that
\begin{enumerate}
\item\label{theo1.1.1} for any $r\leq r_*$, $\Theta_\Gamma(r,\omega_\Gamma(r))$ diverges,
\item\label{theo1.1.2} at $r=r_\sharp$, $\Theta_\Gamma(r,\omega_\Gamma(r))$ converges.
\end{enumerate}
\end{theo}

\begin{rema}
We are unable to determine whether the second assertion is true for any $r>r_*$. If this were true, we would obtain a phase transition for the divergence of the Poincaré series and also for all the assertions in Corollary~\ref{transitionphaseCor} below.
\end{rema}

We are now moving on explaining several applications of Theorem~\ref{mainthm}. Before that, let us put into context the series $\Theta_\Gamma(r,s)$ under consideration. The question whether $\Theta_\Gamma(r,s)$ converges or diverges at $s=\omega_\Gamma(r)$ is of particular importance for many properties such as the construction of boundary measures and growth problems. It is also related to the parabolic gap property coined in~\cite{DWY} that we discuss next.


Consider a relatively hyperbolic group $\Gamma$ and a maximal parabolic subgroup $P$. The \emph{growth rate of the Green function} induced on $P$ is defined as
\[
\omega_P(r)=\limsup_{n\to\infty}\frac{1}{n}\log \sum_{\underset{|x|=n}{x\in P}}G(e,x|r).
\]
Following~\cite{DWY}, we say that $\Gamma$ has a {parabolic gap along $P$ for the Green function} at $r\in[1,R]$ if $\omega_P(r)<\omega_\Gamma(r)$. If this holds along every parabolic subgroups, we say $\Gamma$ has a {parabolic gap for the Green function} at $r$ and if this in turn holds for every $r\in [1,R]$ we say that $\Gamma$ has a \emph{parabolic gap for the Green function}.

This notion is analogous to the parabolic gap condition introduced in~\cite{DOP}, where first examples of convergent (standard) Poincaré series were constructed without this parabolic gap condition. As an interesting consequence, the authors of~\cite{DOP} produced Patterson--Sullivan measures having atoms at parabolic points in the visual boundary of Hadamard manifolds. Convergence or divergence of Poincaré series and the parabolic gap condition are also related to the growth tight\-ness property.


In~\cite{DWY}, we proved that if there exists a relatively hyperbolic group with convergent Poincaré series, then the parabolic gap condition fails. On the other hand, having a parabolic gap has various applications related to asymptotic properties of branching random walks, see in particular~\cite[Theorem~1.8, Remark~5.4]{DWY}.


Motivated by this discussion and by the work of~\cite{DOP}, we make use of the above theorem to further develop various properties for Poincaré series associated with the Green function. We introduce a family of Patterson--Sullivan type measures $\nu_e(r)$ associated with the Poincaré series defined in~\eqref{PoincareSeriesEQ}. We also introduce a family of proper distances $\fkd_r$ which are quasi-isometric to the word distance. Recall that a group $\Gamma$ endowed with a distance $d$ is called growth tight if for every non-trivial quotient $\bar{\Gamma}$ endowed with the corresponding quotient distance $\bar{d}$, the growth rate of $(\bar \Gamma,\bar d)$ is smaller than the growth rate of $(\Gamma,d)$, see Subsection~\ref{s:growthtightness} for more details.


We summarize here the different results we obtain, see Theorem~\ref{transitionphasePoincare}, Corollary~\ref{transitionphasePattersonSullivan} and Theorem~\ref{transitionphasegrowthtightness}.
\begin{coro}\label{transitionphaseCor}
The pair $(\Gamma,\mu)$ in Theorem~\ref{mainthm} has the following properties:
\begin{enumerate}
\item\label{coro1.3.1} the parabolic gap for the Green function holds for $r\le r_*$ but fails at $r=r_\sharp$,
\item\label{coro1.3.2} the Patterson--Sullivan measure $\nu_e(r)$ is supported on conical limit points for $r\le r_*$ and is purely atomic and supported on parabolic limit points at $r=r_\sharp$,
\item\label{coro1.3.3} the growth tightness for the proper distance $\fkd_r$ holds for $r\le r_*$ but fails at $r=r_\sharp$.
\end{enumerate}
\end{coro}


\begin{rema}
In~\cite{Peigne11}, Peigné constructed a divergent Schottky group without the parabolic gap condition. In view of our examples, it is relevant to ask here whether there exist examples of Poincaré series~\eqref{PoincareSeriesEQ} which are divergent, but without a parabolic gap for the Green function.
\end{rema}

\begin{rema}
The proper metric in the third item is defined by
\[
\fkd_r(x,y):=\omega_\Gamma(r)\left|x^{-1}y\right|+\left|x^{-1}y\right|_r
\]
so it is a linear combination of the word metric $|x^{-1}y|$ and the $r$-Green metric $|x^{-1}y|_r:=-\log \frac{G(x,y|r)}{G(e,e|r)}$. It is, in fact, quasi-isometric to a word distance for $r<R$. Our result allows us to answer a variant of a question raised in~\cite{ACT}, namely whether growth tightness is invariant among cocompact actions on geodesic metric spaces. Note that Cashen--Tao~\cite{CaTao16} showed examples of product groups with growth tightness for one generating set but not for another generating set. The above examples within the class of relatively hyperbolic groups are new.
\end{rema}


We now briefly outline the proof of Theorem~\ref{mainthm}. As fore-mentioned, the construction of a convergent Poincaré series was one of the unsolved problems of~\cite{DWY}. However, we did explain there how it might be possible to perform such a construction. Namely, in~\cite[Example~C]{DWY}, we proved the following.

Assume that there exists a finitely generated group $\Gamma_0$ endowed with a probability measure $\mu_0$ such that the associated Poincaré series $\Theta_0(s,r)$ is convergent at $s=\omega_{{\Gamma_0}}(r_0)$ for some fixed $r_0$. Then, it is possible to construct a convergent Poincaré series on the free product $\Gamma_0*\Z^d$. The proof consists in bounding the Poincaré series $\Theta(\omega(r),r)$ on $\Gamma_0*\Z^d$ in terms of the quantity $\sum_{k}\Theta_0(\omega_{{\Gamma_0}}(r_0),r_0)^k\Theta_1(s_1,r_1)^k$, where $\Theta_1(s_1,r_1)$ is a suitable Poincaré series defined on the free factor $\Z^d$. Then, by shifting the weight toward the free factor $\Gamma_0$, it is possible to construct probability measures on $\Gamma_0*\Z^d$ in such a way that $\Theta_1(s_1,r_1)$ is arbitrarily small. As a consequence, if $\Theta_0(\omega_{{\Gamma_0}}(r_0),r_0)$ is finite, then the sum $\sum_{k}\Theta_0(\omega_{{\Gamma_0}}(r_0),r_0)^k\Theta_1(s_1,r_1)^k$ is also finite and so the initial Poincaré series $\Theta(\omega(r),r)$ converges.

In fact, note that the conclusions of Theorem~\ref{mainthm} are a bit stronger. Indeed we do not only construct a convergent series, but we prove that convergence depends on $r$. Therefore, the technicality of~\cite[Example~C]{DWY} needs to be refined a bit. However, the overall strategy remains the same.

In other words, the problem of the construction of a convergent Poincaré series on a relatively hyperbolic group reduces to the problem of finding one such convergent series on some finitely generated group, not necessarily relatively hyperbolic.


Elaborating on the work of Picardello and Woess~\cite{PicardelloWoessbitree}, we prove the following phase transition result for random walks on bi-trees. Consider the Cartesian product of two regular trees $T_1$ and $T_2$ of respective degree $l_1$ and $l_2$. Let $\mu_i$ be {the driving measure for the lazy simple random walk} on $T_i$ and set $\mu=\alpha \mu_1 +(1-\alpha)\mu_2$ for $0<\alpha<1$.

\begin{theo}[{Theorem~\ref{thmbitree}}]\label{mainthmbitree}
The probability measure $\mu$ on $T_1\times T_2$ satisfies the following. If $l_1=l_2$, then for every fixed $r<R$, $H_r(n)\asymp \rme^{n\omega_\Gamma(r)}$ for all $n\ge 1$. If $l_1>l_2$, then there exists a phase transition at some $r_0\in (1,R)$ such that the following holds.
\begin{enumerate}
\item\label{theo1.6.1} For every $r<r_0$, we have $H_r(n)\asymp \rme^{n\omega_\Gamma(r)}$ for all $n\ge 1$.
\item\label{theo1.6.2} At $r=r_0$, we have $H_r(n)\asymp n^{-1}\rme^{n\omega_\Gamma(r)}$ for all $n\ge 1$.
\item\label{theo1.6.3} For every $r_0<r<R$, we have $H_r(n)\asymp n^{-3/2}\rme^{n\omega_\Gamma(r)}$ for all $n\ge 1$.
\end{enumerate}
\end{theo}

As a conclusion, for the group $\Gamma_0=F_2\times F_3$, where $F_k$ is the free groups with $k$ generators, we find that the series $\Theta_0(\omega_{{\Gamma_0}}(r),r)$ diverges for $r\leq r_0$, but converges for $r_0<r<R$. This allows us to prove Theorem~\ref{mainthm}.

Using the example given by Theorem~\ref{mainthmbitree}, we can thus construct a convergent Poincaré series on a relatively hyperbolic group such that the parabolic subgroups are {direct products of free groups of different ranks}. A next interesting problem is to find other classes of parabolic subgroups for which it is possible to construct a convergent Poincaré series, or in the opposite direction, to rule out certain classes of parabolic subgroups.

Beyond free products, standard examples of relatively hyperbolic groups include limit groups appearing as limits of free groups, fundamental groups of finite co-volume hyperbolic manifolds or more generally fundamental groups of geometrically finite Riemannian manifolds of pinched negative curvature. In the two former cases, parabolic subgroups are virtually abelian and in the latter case, they are virtually nilpotent. Therefore, it seems interesting to understand the behavior of $H_r(n)$ for virtually nilpotent groups in view of our initial problem. In a separate paper, we prove that $H_1(n)$ is linear in $n$ for abelian groups and for certain nilpotent groups, as well as other amenable groups. In particular, the associated Poincaré series fails to converge.

In fact, we prove here that whenever parabolic subgroups are amenable, the parabolic gap condition holds and so the Poincaré series must be divergent. Precisely, we prove the following.

\begin{theo}[Corollary~\ref{coramenableparabolic}]
\label{mainthmamenable}
Let $\Gamma$ be a relatively hyperbolic group endowed with a finitely supported admissible and symmetric probability measure $\mu$. Assume that maximal parabolic subgroups of $\Gamma$ are amenable. Then $\Gamma$ has a parabolic gap for the Green function.
\end{theo}

This rules out amenable groups as examples of parabolic subgroups for which there is a convergent Poincaré series, or even a divergent Poincaré series, but without the parabolic gap condition. {Direct} products of free groups are not amenable and so they cannot arise as parabolic subgroups for the aforementioned examples of relatively hyperbolic groups. Note, however, that {direct} products of free groups are right-angled Artin groups. In fact, the group $\Gamma$ in Theorem~\ref{mainthm}, constructed using the example of Theorem~\ref{mainthmbitree}, can be written as $\Gamma=(F_k\times F_l)*\Z^d$ and so it is also a right-angled Artin group. Random walks on such groups have received a lot of attention recently. This gives further motivation for our main theorem.



Let us give more details on how the paper is organized. It has four sections besides the Introduction. Sections~\ref{Sectiongrowthrateat1} to~\ref{Sectionpattersonsullivan} contain preliminary results, whereas Theorem~\ref{mainthm} and Corollary~\ref{transitionphaseCor} are proved in Section~\ref{SectionRHG}.

Specifically, in Section~\ref{Sectiongrowthrateat1}, we consider the growth rate $\omega_\Gamma(r)$ at the special value $r=1$. We prove in particular that $\omega_\Gamma(1)=0$ and study further continuity of the function $r\mapsto \omega_\Gamma(r)$ at $1$. The results of this section will be mainly used to prove Theorem~\ref{mainthmamenable}, that is, the parabolic gap condition must hold and so Poincaré series must be divergent whenever parabolic subgroups are amenable.

Then, in Section~\ref{Sectionbitree}, we prove Theorem~\ref{mainthmbitree}. As a particular consequence, we see that the quantity $H_r(n)$ may fail to be sub-multiplicative. This disproves an argument due to Candellero, Gilch and Müller~\cite{CandelleroGilchMuller} as well as its consequences, see Remark~\ref{remarkCGM} for more details.

In Section~\ref{Sectionpattersonsullivan}, we develop a Patterson--Sullivan theory associated with the Poincaré series $\Theta(s,r)$. This section is not directly related to Theorem~\ref{mainthm}, i.e. the construction itself of a convergent Poincaré series. However, this is the section in which we prove all intermediate results needed to derive Corollary~\ref{transitionphaseCor} from Theorem~\ref{mainthm}. We also recall there the proper definition of relatively hyperbolic subgroups as well as standard results for random walks on such groups.

Finally, in Section~\ref{SectionRHG}, we gather all the results of the previous sections to prove our main theorems. First, we focus on amenable parabolic subgroups and we prove Theorem~\ref{mainthmamenable}. Then, we use Theorem~\ref{mainthmbitree} to construct a convergent Poincaré series and prove Theorem~\ref{mainthm}. Finally, we prove Corollary~\ref{transitionphaseCor}, using the results of Section~\ref{Sectionpattersonsullivan}.

\subsection*{Acknowledgments}
We are grateful to the referee for several helpful comments and suggestions that helped us greatly improve the exposition of the text.



\section{The growth rate of the Green function at 1}
\label{Sectiongrowthrateat1}
We study $\omega_\Gamma(1)$ in this section. We first show that $\omega_\Gamma(r)$ always vanishes at 1 and then study continuity at this special value.



\subsection{Nullity of the growth rate at 1}

In this subsection, we do not need to assume that the random walk is symmetric, nor that it is finitely supported. We start with the following lower bound.

\begin{lemm}\label{lowerboundgrowthrateat1}
Let $\Gamma$ be a finitely generated group and $\mu$ an admissible probability measure on $\Gamma$. Then, $\omega_\Gamma(1)\geq 0$.
\end{lemm}

\begin{proof}
Since the series $\sum_{n = 0}^{\infty} H_1(n) = \sum_{k=0}^{\infty} \sum_{x \in \Gamma} p_k(e, x) = \sum_{k=0}^{\infty} 1 = \infty$, we have that
\begin{equation*}
\omega_{\Gamma}(1) = \limsup_{n \to \infty} \frac{1}{n} \log H_1(n) \geq 0. \qedhere
\end{equation*}
\end{proof}


\begin{rema}
If $\mu$ is finitely supported, then there is a constant $C > 0$ such that $H_1(n) \geq C$ for all $n \geq 1$. In fact, since the random walk driven by {$\mu$ is transient} and has bounded jumps of length at most $C_0>0$, it will eventually visit the annulus
\[
A(n,n+C_0)=\{x,n\leq |x|\leq n+C_0\}.
\]
Note that
\begin{align*}
\sum_{x\in A(n,n+C_0)}G(e,x) &= \sum_{x\in A(n,n+C_0)}\;\sum_{k\ge 0}p_k(e,x)\\
&=\sum_{x\in A(n,n+C_0)}\;\sum_{k\ge 0}\bfP(X_k=x).
\end{align*}
Thus,
\begin{align*}
\sum_{x\in A(n,n+C_0)}G(e,x) & \ \geq \bfP(\text{the random walk ever visits }A(n,n+C_0)) \\
&\ \geq 1.
\end{align*}
Now by~\cite[Lemma~3.1$\MK$(1)]{DWY}, $H_1(n)\asymp H_1(n+1)$, hence
\[
\sum_{x\in A(n,n+C_0)}G(e,x)\asymp H_1(n),
\]
which concludes the proof that $H_1(n) \geq C$.
\end{rema}


\begin{prop}\label{upperboundgrowthrateat1}
Let $\Gamma$ be a finitely generated group and $\mu$ be an admissible probability measure on $\Gamma$. Then $\omega_{\Gamma}(1) \leq 0$. Moreover, if $\Gamma$ has exponential growth, there exists $C>0$ such that $H_1(n)\leq C n^{3}$.
\end{prop}

\begin{proof}
Note that
\[
H_1(n) \leq G(e, e) \sharp S_n,
\]
hence, if the volume growth rate
\[
v = \limsup_{n \to \infty} \frac{1}{n} \log \sharp S_n = 0,
\]
then $\omega_{\Gamma}(1) \leq 0$.

We assume that $v > 0$ in the remainder of the proof. Then there exist $c_1 > 1$ and $v_1$, $v_2 > 0$ such that
\begin{equation}\label{hksubexp}
c_1^{-1} \rme^{v_1 n} \leq \sharp B_n \leq c_1 \rme^{v_2 n},
\end{equation}
where $\sharp B_n$ is the cardinality of the ball of radius $n$. By~\cite[Theorem~1]{V91},
\[
p_n(x, y) \leq c_2 \rme^{- c_3 n^{1/3}}
\]
for some $c_2$, $c_3 > 0$. Choose $c_4$ so that $c_4^{1/3} c_3 > v_2$. Then
\begin{align*}
H_1(n) &\ \leq \sum_{k = 0}^{c_4 n^3}\;\sum_{x \in S_n} p_k(e, x) + c_2\sum_{x \in S_n}\;\sum_{k= c_4 n^3+1}^{\infty} c_3 \rme^{- c_3 k^{1/3}}\\
&\ \leq c_4n^3 + 1 + c_5 \sum_{k = c_4 n^3+1}^{\infty} \rme^{v_2 n} \rme^{- c_3 k^{1/3}}.
\end{align*}


Note that
\[
\sum_{k = c_4 n^3+1}^{\infty} \rme^{- c_3 k^{1/3}} \leq \int_{c_4 n^3}^{\infty} \rme^{-c_3 t^{1/3}} \rmd t \leq c_6 n^{2} \rme^{- c_3 c_4^{1/3} n}.
\]
Thus we have
\[
H_1(n) \leq c_4 n^3 + 1+c_7 n^{2} \rme^{\left(v_2 - c_3 c_4^{1/3} \right) n} \leq c_8 n^3
\]
and hence $\omega_{\Gamma}(1) \leq 0$.
\end{proof}

Lemma~\ref{lowerboundgrowthrateat1} and Proposition~\ref{upperboundgrowthrateat1} together yield the following corollary.
\begin{coro}
Let $\Gamma$ be a finitely generated group and let $\mu$ be an admissible probability measure on $\Gamma$. Then, $\omega_\Gamma(1)=0$.
\end{coro}

We also have the following result.

\begin{prop}\label{omeganonincreasing}
Let $\Gamma$ be a finitely generated group and let $\mu$ be an admissible probability measure on $\Gamma$. Then, $r\mapsto\omega_\Gamma(r)$ is monotonically non-decreasing.
\end{prop}

\begin{proof}
Since the Green function itself is non-decreasing in $r$, given $s\leq r$, we have $H_n(s)\leq H_n(r)$, hence $\omega_\Gamma(s)\leq \omega_\Gamma(r)$.
\end{proof}

\subsection{Continuity of the growth rate}

We start with the following result which holds for every finitely generated group, without assuming that the random walk is finitely supported. We assume however that it is symmetric.

\begin{prop}\label{omegaGcont}
Let $\Gamma$ be a finitely generated group and let $\mu$ be a symmetric admissible probability measure on $\Gamma$. Then, the function $\omega_{\Gamma}(r)$ is continuous for $0 < r < R$.
\end{prop}


\begin{proof}
We modify the arguments in~\cite[Lemma~3.1]{DWY}, where $\mu$ was assumed to be finitely supported and only the case $1<r<R$ was treated. There are constants $c_1 > 0$ and $v_2 \geq 0$ such that $\sharp S_n \leq c_1 \rme^{v_2 n}$. Fix $\delta > 0$. We choose $c_2$ so that for every $\delta \leq r \leq R - \delta$,
\[
v_2 - c_2 \left(\log R - \log (R - r) \right) < \omega_{\Gamma}(r).
\]
Note that since the underlying random walk is symmetric, for every $x$ and every $k$, we have $p_k(e,x)p_k(e,x)\leq p_{2k}(e,e)$ and by~\cite[Lemma~1.9]{Woessbook}, $p_{2k}(e,e)\leq R^{-2k}$. Thus,
\begin{equation}\label{decayheatkernel}
p_k(e, x) \leq R^{-k}
\end{equation}
for every $x \in \Gamma$ and $k \ge 0$. Consequently, we have for $\delta \le r \le R - \delta$,
\[
\sum_{x \in S_n}\;\sum_{k\,>\,c_2 n} r^k p_k(e, x) \le c_1 \rme^{v_2 n} \sum_{k\,>\,c_2 n} \left(\frac{r}{R} \right)^k \le c_1 \delta^{-1} R \rme^{v_2 n - c_2 \left(\log R - \log r \right) n}.
\]
By the choice of $c_2$ we have that
\[
\omega_{\Gamma}(r) = \limsup_{n \to \infty} \frac{1}{n} \log \sum_{x \in S_n}\;\sum_{k=0}^{c_2 n} r^k p_k(e, x).
\]
Now for $0 < s < r \le R - \delta$,
\[
H_n(s) \geq \sum_{x \in S_n}\;\sum_{k = 0}^{c_2 n} s^k p_k(e, x)
\ge \left(\frac{s}{r} \right)^{c_2 n} \sum_{x \in S_n}\;\sum_{k = 0}^{c_2 n} r^k p_k(e, x),
\]
and hence $0 \leq \omega_{\Gamma}(r) - \omega_\Gamma(s) \le c_2 (\log r - \log s )$, since $\omega_\Gamma$ is non-decreasing by Proposition~\ref{omeganonincreasing}. Now $\delta > 0$ is arbitrary, so this proves that $\omega_{\Gamma}(r)$ is continuous in $0 < r < R$.
\end{proof}



The continuity at the inverse of the spectral radius seems to be a difficult problem in general. It is already known that $\omega_\Gamma$ is continuous at $R$ in hyperbolic groups~\cite[Theorem~1.1]{SWX} and in relatively hyperbolic groups~\cite[Theorem~1.1]{DWY}. However, we do not know much beyond these classes of groups.


We now investigate further continuity at $r=1$. By~\cite[Theorem~12.5]{Woessbook}, $R=1$ can only happen if $\Gamma$ is amenable, hence the following discussion only applies to amenable groups.

Assume that the random walk satisfies the following Gaussian lower bound. There exists a sub-exponential function $f$, i.e. $\frac{1}{n}\log f(n)\to 0$, $n\to \infty$, such that for every~$k$, for every $x\in \Gamma$,
\begin{equation}\label{Gaussianlowerbound}
p_k(e, x) \gtrsim f(k) \rme^{- c \frac{|x|^2}{k}}.
\end{equation}
Then,
\begin{align*}
H_r(n) =& \sum_{x \in S_n}\;\sum_{k = 0}^{\infty} r^k p_k(e, x)
\gtrsim \sharp S_n \sum_{k = 0} f(k) \rme^{- k \log r^{-1} - c \frac{n^2}{k}} \\
\gtrsim & \ f\left(\left \lfloor \frac{n}{\sqrt{\log r^{-1}}} \right \rfloor \right)\sharp S_n \rme^{- n (1 + c) \sqrt{\log r^{-1}}}.
\end{align*}
It follows that
\[
\omega_{\Gamma}(r) \geq -(1 + c) \sqrt{\log r^{-1}} + v.
\]
Letting $r \uparrow 1$, we see that $0 \geq \limsup_{r \uparrow 1} \omega_{\Gamma}(r) \geq v$. Therefore $v = 0$ and
\[
\lim_{r \uparrow 1} \omega_{\Gamma}(r) = 0.
\]



\begin{rema}
We thus recover the fact that a Gaussian lower bound like~\eqref{Gaussianlowerbound} cannot hold for groups of exponential volume growth, which was already noticed in literature, see for instance the comments after~\cite[(0.3)]{Alexopoulospolycyclic}.
\end{rema}

\begin{prop}
Let $\Gamma$ be a finitely generated virtually nilpotent group endowed with a finite generating set and let $\mu$ be a finitely supported symmetric admissible probability measure on $\Gamma$. Then, the function $\omega_\Gamma$ is left-continuous at 1.
\end{prop}

\begin{proof}
The fact that a Gaussian lower bound~\eqref{Gaussianlowerbound} with $f(k)=k^{-d/2}$ holds for virtually nilpotent groups is well known, see for instance~\cite[Corollary~1.9]{Alexopoulos}.
\end{proof}


\section{Cartesian products of trees}\label{Sectionbitree}
In this section, we study the behavior of $H_r(n)$ for random walks on $T\times T'$, where $T,T'$ are regular trees. The asymptotics for the Green functions are given by the work of Picardello and Woess~\cite{PicardelloWoessbitree}, see also~\cite[Section~28]{Woessbook} and references therein for the particular case of $T\times \Z$.

Let $T_1,T_2$ be regular trees of degree $l_1,l_2\geq 3$. We consider the lazy simple random walk $\mu_i$ on $T_i$ whose transition kernel $p_i(x,y)$ is defined by
\[
p_i(x,y)=
\begin{cases}
\frac{1}{2l_i} &\text{if } x,y \text{ are connected with an edge in } T_i,\\
\frac{1}{2} &\text{if }x=y,\\
0 &\text{otherwise.}
\end{cases}
\]
In particular, $\mu_i$ is an admissible symmetric finitely supported probability measures on $T_i$. For every $\alpha_1,\alpha_2\geq 0$, $\alpha_1+\alpha_2=1$, we let $\mu$ be the probability measure on $T_1\times T_2$ given by
\begin{equation}\label{defmubitree}
\mu=\alpha_1 \mu_1+\alpha_2\mu_2.
\end{equation}
In terms of Markov operators, this means that
\[
P_{\mu}=\alpha_1 P_{\mu_1}\otimes I+\alpha_2 I\otimes P_{\mu_2}.
\]
As noted in~\cite[Section~3]{PicardelloWoessbitree}, the lazy simple random walk on $T_1\times T_2$ with $\mu(x,x)=1/2$ is given by $\alpha_i=\frac{l_i}{l_1+l_2}$. Set $\rho_i = \frac{1}{2} + \frac{\sqrt{l_i - 1}}{l_i}$ and $R = \frac{1}{\alpha_1 \rho_1 + \alpha_2 \rho_2}$. Then, $\rho_i$~is the spectral radius of $\mu_i$, see~\cite[(2.2)]{PicardelloWoessbitree} and $\rho=R^{-1}$ is the spectral radius of $\mu$, see~\cite[Section~3]{PicardelloWoessbitree}. We prove here the following.
\begin{theo}\label{thmbitree}
If $l_1=l_2$, then for every $r<R$, we have $H_r(n)\asymp \rme^{n\omega_\Gamma(r)}$. If $l_1>l_2$, then there exists a phase transition at some $r_0\in (1,R)$ such that the following holds.
\begin{itemize}
\item For every $r<r_0$, we have $H_r(n)\asymp \rme^{n\omega_\Gamma(r)}$.
\item At $r=r_0$, we have $H_r(n)\asymp n^{-1}\rme^{n\omega_\Gamma(r)}$.
\item For every $r_0<r<R$, we have $H_r(n)\asymp n^{-3/2}\rme^{n\omega_\Gamma(r)}$.
\end{itemize}
\end{theo}


\begin{proof}
The remainder of the section is devoted to the proof of this theorem. By~\cite[Theorem~3.1]{PicardelloWoessbitree}, for every $r<R$, for every $\lambda_0\geq 0$, there exist $r_1,r_2$ such that as $x=(x_1,x_2) \in T_1\times T_2$ tends to infinity and $|x_2|/|x_1|=\lambda$ converges $\lambda_0$, we have
\[
G(e,x|r)\sim G_1(e,x_1|r_1)G_2(e,x_2|r_2)\bigg(|x_1|+\frac{l_1}{l_1-2}\bigg)\bigg(|x_2|+\frac{l_2}{l_2-2}\bigg)\frac{1}{|x_1|^{5/2}}C(\lambda),
\]
where $C(\lambda)$ is a continuous positive function.

The numbers $r_1$ and $r_2$ are the unique solutions of the system
\begin{equation}\label{systemr1r2}
\begin{cases}
\alpha_1r_1^{-1}+\alpha_2r_2^{-1}=r^{-1}\\
\alpha_2\sqrt{\left(r_2^{-1}-\frac{1}{2}\right)^2-\frac{l_2-1}{l_2^2}}=\lambda \alpha_1\sqrt{\left(r_1^{-1}-\frac{1}{2}\right)^2-\frac{l_1-1}{l_1^2}}
\end{cases}
\end{equation}
By symmetry, the same holds when $|x_1|/|x_2|=\lambda'$ converges to $\lambda'_0\in [0,+\infty)$, switching the indices $1$ and $2$ and replacing the function $C$ with a function $C'$ which is also continuous and positive.

Set $\beta_i = \frac{\sqrt{l_i - 1}}{l_i}$ and
\begin{align*}
F_i(r) &= \frac{l_i}{l_i - 1} \left(r^{-1} - \frac{1}{2} - \sqrt{\left(r^{-1} - \frac{1}{2} \right)^2 - \beta_i^2} \right), \\
G_i(r) &= \frac{r^{-1}}{r^{-1} - (1/2) \left(1 + F_i(r) \right)}.
\end{align*}
By~\cite[(2.4)]{PicardelloWoessbitree}, we have that
\[
G_i(e_i, \, x_i | r) = G_i(r) F_i(r)^{|x_i|},
\]
hence in particular $G_i(e_i,e_i|r)=G_i(r)$.

Since $1 \leq G_i(r) \leq G_i(\rho_i^{-1})$, we have
\begin{equation}\label{Hniasymp}
H_n^i(r) = \sum_{x_i \in T_i \colon |x_i| =n} G_i(e_i, x_i | r) = l_i (l_i - 1)^{n - 1} G_i(r) F_i(r)^n \asymp \rme^{n \omega_{T_i}(r)},
\end{equation}
where
\[
\omega_{T_i}(r) = \log (l_i - 1) + \log F_i(r).
\]

We now fix $r<R$ and we consider $n\geq0$. The sphere $S_n$ in $T_1\times T_2$ can be decomposed as
\[
S_n=\bigcup_{k=0}^nS_{k}^1\times S_{n-k}^2.
\]
Then, for every $k\leq n$, there exists a unique couple $(r_1(\lambda), r_2(\lambda))$ satisfying~\eqref{systemr1r2} with $\lambda=\frac{n-k}{k}$. Moreover,
\[
\frac{1}{|x_1|^{5/2}}C(\lambda)=\frac{1}{(|x_1|+|x_2|)^{5/2}}(1+\lambda)^{5/2}C(\lambda)
\]
and for $k\geq n/2$, $\lambda \leq 1$, hence by the continuity of $C(\lambda)$,
\[
\frac{1}{|x_1|^{5/2}}C(\lambda)\asymp \frac{1}{(|x_1|+|x_2|)^{5/2}}.
\]
Similarly,
\[
\frac{1}{|x_2|^{5/2}}C'(\lambda')=\frac{1}{(|x_1|+|x_2|)^{5/2}}(1+\lambda')^{5/2}C'(\lambda')
\]
and so for $k\leq n/2$,
\[
\frac{1}{|x_2|^{5/2}}C'(\lambda')\asymp \frac{1}{(|x_1|+|x_2|)^{5/2}}.
\]
Setting $\kappa_i=\frac{l_i}{l_i-2}$, we thus have
\begin{multline*}
\sum_{x\in S_n}G(e,x|r)\\
\begin{aligned}
&\asymp \frac{1}{n^{5/2}}\sum_{k=0}^n \;
\sum_{x_1\in S_k^1}\;\sum_{x_2\in S_{n-k}^2}G_1(e,x_1|r_1(\lambda)) G_2(e,x_2|r_2(\lambda))(k+\kappa_1)((n-k)+\kappa_2)\\
&\asymp \frac{1}{n^{5/2}}\sum_{k=0}^n(k+\kappa_1)((n-k)+\kappa_2) H^1_k(r_1(\lambda)) H^2_{n-k}(r_2(\lambda)).
\end{aligned}
\end{multline*}
Applying~\eqref{Hniasymp}, we see that
\begin{equation}\label{Hnasymp}
H_n(r) \asymp \frac{1}{n^{5/2}} \sum_{k = 0}^n (k+\kappa_1)((n-k)+\kappa_2) \exp \left(n \Psi(\lambda) \right),
\end{equation}
with $\lambda = \frac{n - k}{k}$ and
\[
\Psi(\lambda) = \frac{1}{1 + \lambda} \omega_{T_1}(r_1(\lambda)) + \frac{\lambda}{1 + \lambda} \omega_{T_2}(r_2(\lambda)).
\]
In order to find the asymptotics of $H_n(r)$, we thus need to find where the function $\Psi(\lambda)$ takes its maximal value.


Now let $t = r^{-1} - \frac{1}{2}$ and $t_i = t_i(\lambda) = r_i(\lambda)^{-1} - \frac{1}{2}$. Then $(t_1, t_2)$ solves the system of equations
\begin{equation}\label{systemt1t2}
\begin{cases}
\alpha_1 t_1 + \alpha_2 t_2 = t, \\
\alpha_2 \sqrt{t_2^2 - \beta_2^2} = \lambda \alpha_1 \sqrt{t_1^2 - \beta_1^2}.
\end{cases}
\end{equation}

\begin{lemm}\label{t1t2deriviative}
The functions $\lambda \mapsto t_1(\lambda)$ and $\lambda \mapsto t_2(\lambda)$ are continuously differentiable. Furthermore,
\begin{align*}
t_1'(\lambda) &= - \frac{\lambda \alpha_1\left(t_1^2 - \beta_1^2\right)}{\alpha_2 t_2 + \lambda^2 \alpha_1 t_1}, \\
t_2'(\lambda) &= - \frac{\alpha_1}{\alpha_2} t_1'(\lambda) = \frac{\lambda^{-1} \alpha_2 \left(t_2^2 - \beta_2^2 \right)}{\alpha_2 t_2 + \lambda^2 \alpha_1 t_1}.
\end{align*}
\end{lemm}

\begin{proof}
Let $U$ be the open set $(\beta_1,+\infty)\times (\beta_2,+\infty)\times (\beta,+\infty)\times (0,+\infty)$ with $\beta = \rho - \frac{1}{2}$. We set
\[
\Upsilon:(t_1,t_2,t,\lambda)\in U\mapsto \left(\alpha_1 t_1 + \alpha_2 t_2-2t, \alpha_2 \sqrt{t_2^2 - \beta_2^2}- \lambda \alpha_1 \sqrt{t_1^2 - \beta_1^2}\right).
\]
Then
\begin{align*}
\frac{\partial \Upsilon}{\partial t}&=(-2,0)
\\
\intertext{and}
\frac{\partial \Upsilon}{\partial \lambda}&=\left(0,- \alpha_1 \sqrt{t_1^2-\beta^2}\right).
\end{align*}
For $t_1>\beta_1$, the matrix
\[
\begin{pmatrix}
-2&0\\
0&- \alpha_1 \sqrt{t_1^2-\beta^2}
\end{pmatrix}
\]
is invertible. The implicit function theorem shows that the solution $(t_1,t_2)$ of~\eqref{systemt1t2} is continuously differentiable in the variables $(t,\lambda)$. The formulas for $t_1'(\lambda)$ and $t_2'(\lambda)$ are then derived from~\eqref{systemt1t2}.
\end{proof}

Define
\[
\varphi_i(t) = \log l_i + \log \left(t - \sqrt{t^2 - \beta_i^2} \right).
\]
Then
\[
\Psi(\lambda) = \frac{1}{1 + \lambda} \varphi_1(t_1(\lambda)) + \frac{\lambda}{1 + \lambda} \varphi_2(t_2(\lambda)).
\]
Since
\begin{equation}\label{phideriviative}
\varphi_i'(t) = - \frac{1}{\sqrt{t^2 - \beta_i^2}},
\end{equation}
we have that
\[
\lambda \varphi_2'(t_2(\lambda)) = \frac{\alpha_2}{\alpha_1} \varphi_1'(t_1(\lambda)),
\]
and hence
\begin{equation}\label{Psideriviative}
\begin{split}
\Psi'(\lambda) &= \frac{\varphi_1'(t_1(\lambda))}{1 + \lambda} t_1'(\lambda) + \frac{\lambda \varphi_2'(t_2(\lambda))}{1 + \lambda} t_2'(\lambda) + \frac{\varphi_2(t_2(\lambda)) - \varphi_1(t_1(\lambda))}{(1 + \lambda)^2} \\
&= \frac{\alpha_1 t_1'(\lambda) + \alpha_2 t_2'(\lambda)}{\alpha_1 (1 + \lambda)} \varphi_1'(t_1(\lambda)) + \frac{\varphi_2(t_2(\lambda)) - \varphi_1(t_1(\lambda))}{(1 + \lambda)^2} \\
&= \frac{\varphi_2(t_2(\lambda)) - \varphi_1(t_1(\lambda))}{(1 + \lambda)^2}.
\end{split}
\end{equation}
Furthermore,
\begin{equation}\label{Psi2deriviative}
\Psi''(\lambda) = - \frac{2}{(1 + \lambda)^3} \bigl[\varphi_2(t_2(\lambda)) - \varphi_1(t_1(\lambda)) \bigr] - \frac{\alpha_1 \sqrt{t_1^2 - \beta_1^2}}{(1 + \lambda) \left(\alpha_2 t_2 + \lambda^2 \alpha_1 t_1 \right)}.
\end{equation}

By Lemma~\ref{t1t2deriviative} and~\eqref{phideriviative}, we see that $\varphi_1(t_1(\lambda))$ (resp. $\varphi_2(t_2(\lambda))$) is strictly increasing (resp. decreasing) in $\lambda$. It follows that there is at most one $\lambda_0 \in [0, \, + \infty)$ such that $\Psi'(\lambda_0) = 0$. Note that $t_2(0) = \beta_2$, $t_1(0) = \alpha_1^{-1} \left(t - \alpha_2 \beta_2 \right) > \beta_1$, and
\[
\varphi_2(t_2(0)) = \frac{1}{2} \log \left(l_2 - 1 \right).
\]
On the other hand, $\lim_{\lambda \to + \infty} t_1(\lambda) = \beta_1$, $\lim_{\lambda \to + \infty} t_2(\lambda) = \alpha_2^{-1}(t - \alpha_1 \beta_1) > \beta_2$, and
\[
\lim_{\lambda \to + \infty} \varphi_1(t_1(\lambda)) = \frac{1}{2} \log \left(l_1 - 1 \right).
\]

Assume $l_1 = l_2 = l$. Then for $\lambda_0 = \frac{\alpha_2}{\alpha_1}$ we have that $t_1(\lambda_0) = t_2(\lambda_0) = t$ and $\Psi'(\lambda_0) = 0$. Note that
\[
\Psi''(\lambda_0) = - \frac{\alpha_1^2 \sqrt{t^2 - \beta^2}}{\alpha_2 t} < 0
\]
by Lemma~\ref{t1t2deriviative} and~\eqref{phideriviative}. By Lemma~\ref{asymptotics}$\MK$\eqref{lemm3.3.2} below, we can deduce that
\begin{equation}\label{Hnl1=l2}
H_n(r) \asymp \rme^{n \Psi(\lambda_0)} = \rme^{n \log \left[l \left(r^{-1} - \frac{1}{2} - \sqrt{\left(r^{-1} - \frac{1}{2} \right)^2 - \beta^2} \right) \right]}.
\end{equation}
This concludes the proof of Theorem~\ref{thmbitree} for the case $l_1=l_2$.


Assume now that $l_1 > l_2$. Then, we have that $\beta_1 < \beta_2$ and $\varphi_1(\beta_1) > \varphi_2(\beta_2)$. Since $\lim_{s \to + \infty} \varphi_1(s) = - \infty$, there exists $t_0 > \alpha_1 \beta_1 + \alpha_2 \beta_2 > \beta_1$ such that
\[
\varphi_1 \left(\alpha_1^{-1} \left(t_0 - \alpha_2 \beta_2 \right) \right) = \varphi_2(\beta_2).
\]

If $t < t_0$, then
\[
\varphi_1(t_1(0)) = \varphi_1 \left(\alpha_1^{-1} \left(t - \alpha_2 \beta_2 \right) \right) > \varphi_1 \left(\alpha_1^{-1} \left(t_0 - \alpha_2 \beta_2 \right) \right) = \varphi_2 (t_2(0))
\]
and hence $\Psi'(\lambda) < 0$ for all $\lambda \geq 0$. It follows that $\Psi(\lambda)$ takes its unique maximum at $\lambda = 0$. By Lemma~\ref{asymptotics}$\MK$\eqref{lemm3.3.1} below,
\begin{equation}\label{Hnt<t0}
H_n(r) \asymp n^{-3/2} \rme^{n \log \varphi_1 \left(\alpha_1^{-1} \left(t - \alpha_2 \beta_2 \right) \right)}.
\end{equation}

Similarly, if $t = t_0$, then $0$ is also the unique maximum point of $\Psi(\lambda)$, and we have further that $\Psi'(0) = 0$,
\[
\Psi''(0) = - \frac{\sqrt{\left(t_0 - \alpha_2 \beta_2 \right)^2 - \alpha_1 \beta_1^2}}{\alpha_2 \beta_2} < 0.
\]
Thus, Lemma~\ref{asymptotics}$\MK$\eqref{lemm3.3.3} below shows that
\begin{equation}\label{Hnt=t0}
H_r(n) \asymp n^{-1} \rme^{n \Psi(0)} = n^{-1} \rme^{n \log \varphi_1 \left(\alpha_1^{-1} \left(t - \alpha_2 \beta_2 \right) \right)}.
\end{equation}

It remains to consider the case that $t > t_0$. Since $t_0 > \alpha_1 \beta_1 + \alpha_2 \beta_2$ we have
\[
\lim_{\lambda \to + \infty} \varphi_2 \left(t_2 (\lambda) \right) = \varphi_2 \left(\alpha_2^{-1} \left(t - \alpha_1 \beta_1 \right) \right) < \varphi_2(\beta_2) < \varphi_1(\beta_1) = \lim_{\lambda \to + \infty} \varphi_1(t_1(\lambda)),
\]
and
\[
\varphi_1(t_1(0)) = \varphi_1 \left(\alpha_1^{-1} \left(t - \alpha_2 \beta_2 \right) \right) < \varphi_1 \left(\alpha_1^{-1} \left(t_0 - \alpha_2 \beta_2 \right) \right) = \varphi_2(t_2 (0)).
\]
Thus there exists $\lambda_0 > 0$ such that $\varphi_1(t_1(\lambda_0)) = \varphi_2(t_2(\lambda_0))$. By~\eqref{Psi2deriviative}, $\Psi''(\lambda_0) < 0$, and we see from Lemma~\ref{asymptotics}$\MK$\eqref{lemm3.3.2} below that
\begin{equation}\label{Hnt>t0}
H_r(n) \asymp \rme^{n \varphi_1(t_1(\lambda_0))}.
\end{equation}

To conclude, what is left to do is proving that $r_0>1$, i.e. $t_0<1/2$. Lengthily computations would prove that
\[
\varphi_1 \left(\alpha_1^{-1} \left(\frac{1}{2} - \alpha_2 \beta_2 \right) \right) < \varphi_2(\beta_2),
\]
hence we necessarily have $t_0<1/2$.
However, we see that at $t_0$, we have by~\eqref{Hnt=t0} $ H_r(n) \asymp n^{-1} \rme^{n \Psi(0)} = n^{-1} \rme^{n \log \varphi_1 \left(\alpha_1^{-1} \left(t - \alpha_2 \beta_2 \right) \right)}. $ In particular,
\[
\omega_\Gamma(r_0)=\log \varphi_1 \left(\alpha_1^{-1} \left(t_0 - \alpha_2 \beta_2 \right) \right)=\log \varphi_2(\beta_2)>0.
\]
Since $\omega_\Gamma(1)=0$ and $\omega_\Gamma$ is increasing, we see directly that $r_0>1$. This ends the proof of Theorem~\ref{thmbitree}. \end{proof}



\begin{lemm}\label{asymptotics}
Assume that $\Phi \in C^2 \left([0, +\infty) \right)$ is eventually decreasing and has a unique maximum point at $0 \leq \lambda_0 < \infty$. Define
\[
f(n) = \sum_{k = 0}^n k (n - k) \rme^{n \Phi \left(\!\frac{n - k}{k}\!\right)}.
\]
\begin{enumerate}\romanenumi
\item\label{lemm3.3.1} If $\lambda_0 = 0$ and $\Phi'(0) < 0$, then
\[
f(n) \asymp n \rme^{n \Phi(0)}.
\]
\item\label{lemm3.3.2} If $\lambda_0 > 0$ and $\Phi''(\lambda_0) < 0$, then
\[
f(n) \asymp n^{5/2} \rme^{n \Phi(\lambda_0)}.
\]
\item\label{lemm3.3.3} If $\lambda_0 = 0$, $\Phi'(0) = 0$ and $\Phi''(0) < 0$, then
\[
f(n) \asymp n^{3/2} \rme^{n \Phi(0)}.
\]
\end{enumerate}
\end{lemm}

\begin{proofc}
\begin{proof}[{\eqref{lemm3.3.1}}]
For every $0 < \varepsilon < - \Phi'(0)$, there exists $\delta > 0$ such that for every $0 \leq \lambda \leq \delta$, $\left| \Phi'(\lambda) - \Phi'(0) \right| < \varepsilon$. By the mean value theorem, for $\lambda \leq \delta$,
\[
\left(\Phi'(0) - \varepsilon \right) \lambda \leq \Phi(\lambda) - \Phi(0) \leq \left(\Phi'(0) + \varepsilon \right) \lambda.
\]
If $\frac{n - k}{k} \leq \delta$, then we have $k \geq (1+\delta)^{-1}n$. Thus
\begin{align*}
f(n) &\geq \sum_{k \ge (1+\delta)^{-1} n} k (n - k) \rme^{n \Phi \left(\!\frac{n - k}{k}\!\right)} \\
&\succ n \rme^{n \Phi(0)} \sum_{k \ge (1+\delta)^{-1} n} (n - k) \rme^{n \left(\Phi'(0) - \varepsilon \right) \frac{n - k}{k}} \\
&\geq n \rme^{n \Phi(0)} \sum_{k \le \frac{\delta}{1 + \delta} n} k \rme^{\left(\Phi'(0) - \varepsilon \right) k} \\
&\succ  n \rme^{n \Phi(0)}.
\end{align*}
Since $0$ is the unique maximum point of $\Phi(\lambda)$, there exists $\eta > 0$ such that $\Phi(\lambda) \leq \Phi(0) - \eta$ for all $\lambda \geq \delta$. Thus
\[
\sum_{k\,<\,(1 + \delta)^{-1}n } k (n - k) \rme^{n \Phi \left(\!\frac{n - k}{k}\!\right)}
\leq n^3 \rme^{n \left(\Phi(0) - \eta \right)}.
\]
Also,
\begin{align*}
\sum_{k \ge (1 + \delta)^{-1} n} k (n - k) \rme^{n \Phi \left(\!\frac{n - k}{k}\!\right)}
&\leq n \rme^{n \Phi(0)} \sum_{k \ge (1 + \delta)^{-1} n} (n - k) \rme^{n \left(\Phi'(0) + \varepsilon \right) \frac{n - k}{k}} \\
&\leq n \rme^{n \Phi(0)} \sum_{k \le \frac{\delta}{1 + \delta}n} k \rme^{\left(\Phi'(0) + \varepsilon \right) (1 + \delta)^{-1} k} \\
&\leq n \rme^{n \Phi(0)} \sum_{k = 0}^{\infty} k \rme^{\left(\Phi'(0) + \varepsilon \right) (1 + \delta)^{-1} k}
\end{align*}
Combining the last two displays yields the desired upper-bound.
\let\qed\relax
\end{proof}

\begin{proof}[{\eqref{lemm3.3.2}}]
For $0 < c_1 < - \frac{\Phi''(\lambda_0)}{2} < c_2$, there exists $\delta > 0$ such that
\begin{equation}\label{secondorderapproximationPhi}
- c_2 \left(\lambda - \lambda_0 \right)^2 \leq \Phi(\lambda) - \Phi(\lambda_0) \leq - c_1 \left(\lambda - \lambda_0 \right)^2
\end{equation}
for $\left| \lambda - \lambda_0 \right| \leq \delta$. Clearly, we have that
\[
\sum_{k \colon \left\lvert \frac{n - k}{k} - \lambda_0 \right\rvert\,>\,\delta} k (n - k) \rme^{n \Phi \left(\!\frac{n - k}{k}\!\right)} \prec n^3 \rme^{n \left(\Phi(\lambda_0) - \eta \right)}
\]
for some $\eta > 0$. Now,
\begin{align*}
\sum_{k \colon \left\lvert\frac{n - k}{k} - \lambda_0 \right\rvert \le \delta} k (n - k) \rme^{n \Phi \left(\!\frac{n - k}{k}\!\right)}
&\prec n^3 \rme^{n \Phi(\lambda_0)} \sum_{k \colon \left\lvert \frac{n - k}{k} - \lambda_0 \right\rvert \le \delta} \frac{1}{n} \rme^{- c_1 n \left(\frac{n}{k} - 1 - \lambda_0 \right)^2} \\
&\asymp n^3 \rme^{n \Phi(\lambda_0)} \int_{(1 + \lambda_0 + \delta)^{-1}}^{(1 + \lambda_0 - \delta)^{-1}} \rme^{- c_1 n \left(x^{-1} - 1 - \lambda_0 \right)^2} \rmd x.
\end{align*}
By a change of variables, we get
\begin{align*}
\sum_{k \colon \left\lvert \frac{n - k}{k} - \lambda_0 \right\rvert \le \delta} k (n - k) \rme^{n \Phi \left(\!\frac{n - k}{k}\!\right)}
&= n^3 \rme^{n \Phi(\lambda_0)} \int_{- \delta}^{\delta} \rme^{-c_1 n y^2} \frac{\rmd y}{(y + 1 + \lambda_0)^2} \\
&\asymp n^{5/2} \rme^{n \Phi(\lambda_0)} \int_{- \delta \sqrt{n}}^{\delta \sqrt{n}} \rme^{- c_1 z^2} \rmd z \\
&\asymp n^{ 5 / 2} \rme^{n \Phi(\lambda_0)}.
\end{align*}
The lower bound can be proved by the same arguments, changing $c_1$ with $c_2$.\let\qed\relax
\end{proof}

\begin{proof}[{\eqref{lemm3.3.3}}]
The proof of~\eqref{lemm3.3.3} is similar to that of~\eqref{lemm3.3.2}, except that we need to replace~\eqref{secondorderapproximationPhi} with
\[
- c_2 \lambda^2 \leq \Phi(\lambda) - \Phi(0) \leq - c_1 \lambda ^2
\]
and the order of magnitude is not $\frac{n-k}{k}\asymp \lambda_0\pm \delta$ anymore, but $k\geq (1+\delta)^{-1}n$ as in~\eqref{lemm3.3.1}. \qedhere
\end{proof}\let\qed\relax
\end{proofc}



\section{Twisted Patterson--Sullivan measures}
\label{Sectionpattersonsullivan}

In the next section, we will give applications of our results to relatively hyperbolic groups. Before that, we introduce here all the necessary material to derive Corollary~\ref{transitionphaseCor} from Theorem~\ref{mainthm}. In particular, we introduce a proper distance $\fkd_r$ which is quasi-isometric to the word distance and a family of Patterson--Sullivan type measures on the Bowditch boundary of a relatively hyperbolic group.


\subsection{Some background on relatively hyperbolic groups}
\label{subsectiondefRHG}

Since their introduction by Gromov, relatively hyperbolic groups were studied by many authors through several equivalent definitions. We will mainly use the viewpoint of Bowditch~\cite{Bowditch} and Gerasimov--Potyagailo~\cite{GePoJEMS,GePoGGD,GP16} in the sequel. Consider a finitely generated group $\Gamma$ acting properly via isometries on a proper geodesic Gromov hyperbolic space~$X$. Define the \emph{limit set} $\Lambda_\Gamma$ as the closure of~$\Gamma$ in the Gromov boundary $\partial X$ of $X$, that is, fixing a base point $x_0$ in $X$, $\Lambda_\Gamma$ is the set of all possible limits of sequences $g_n\cdot x_0$ in $\partial X$, $g_n\in \Gamma$. The proper action of $\Gamma$ on~$X$ by isometries extends to a convergence group action on $\Lambda_\Gamma$ by homeomorphisms, which means that the induced action on the space of distinct triples is properly continuous (see~\cite[Proposition~1.11]{Bo99} for example). If $\Lambda_\Gamma$ contains at least three points, then $\Gamma$ acts minimally on $\Lambda_\Gamma$.

A loxodromic element $x\in \Gamma$ is an infinite order element with exactly two fixed points $x_-\ne x_+$ in $\Lambda_\Gamma$. Moreover, $x$ acts via North-South dynamics on $\Lambda_\Gamma$ in the sense that for any $\xi\neq x_{\pm}$, $x^{\mp n}\xi$ converges to $x_{\mp}$ as $n$ goes to $+\infty$. Then $x_+$ is called the attracting fixed point of $x$ and $x_-$ is called its repelling fixed point. The fixed points of any two loxodromic elements are either the same or disjoint. So $\Gamma$ contains infinitely many loxodromic elements with pairwise disjoint fixed points.


A point $\xi\in \Lambda_\Gamma$ is called \emph{conical} if there is a sequence $g_{n}$ of $\Gamma$ and distinct points $\xi_1,\xi_2$ in $\Lambda_\Gamma$ such that $g_{n}\xi$ converges to $\xi_1$ and $g_{n}\zeta$ converges to $\xi_2$ for all $\zeta\neq \xi$ in~$\Lambda_\Gamma$. A point $\xi\in \Lambda_\Gamma$ is called \emph{parabolic} if its stabilizer in $\Gamma$ is infinite, fixes exactly $\xi$ in $\Lambda_\Gamma$ and contains no loxodromic element. If, in addition, its stabilizer in $\Gamma$ acts co-compactly on $\Lambda_\Gamma \setminus \{\xi\}$, then $\xi$ is called \emph{bounded parabolic}. Say that the action of~$\Gamma$ on $X$ is \emph{geometrically finite} if the induced convergence group action on the limit set $\Lambda_
\Gamma$ is geometrically finite: $\Lambda_
\Gamma$ only consists of conical limit points and bounded parabolic limit points. See~\cite{Bo99} for more general facts on convergence groups.


A group $\Gamma$ is called \emph{relatively hyperbolic} with respect to a collection of subgroups~$\bbP$ if it acts geometrically finitely on a proper geodesic hyperbolic space $X$ such that the stabilizers of parabolic limit points are exactly the conjugates of the elements of~$\bbP$. Elements of $\bbP$ are called maximal parabolic subgroups. We will write $\bbP_0$ for the choice of a set of representatives of conjugacy classes of elements of $\bbP$. By~\cite[Proposition~6.15]{Bowditch}, such a set $\bbP_0$ is finite.

The limit set $\Lambda_\Gamma$ in the Gromov boundary of $X$ is called the \emph{Bowditch boundary} of $\Gamma$. By~\cite[Theorem~9.4]{Bowditch}, it is unique up to equivariant homeomorphism and in particular does not depend on the choice of a proper geodesic Gromov hyperbolic space $X$ on which $\Gamma$ acts geometrically finitely. We will write $\partial\Gamma$ for the Bowditch boundary of $\Gamma$ in the sequel. A relatively hyperbolic group is called \emph{non-elementary} if its Bowditch boundary is infinite; equivalently, $\sharp \partial\Gamma>2$.


We now fix a finite set $\bbP_0$ of representatives of conjugacy classes of maximal parabolic subgroups. When $P\in \bbP_0$ and $g\in \Gamma$, we call $gP$ a maximal parabolic coset.
\begin{defi}
Let $gP$ be a maximal parabolic coset and $\eta, L>0$ be fixed constants. A point $p$ on a geodesic $\alpha$ is called \emph{$(\eta, L)$-deep} in $gP$ if
\[
B(p, 2L)\cap \alpha\subseteq N_\eta(gP).
\]
It is called an \emph{$(\eta, L)$-transition point} if it is not \emph{$(\eta, L)$-deep} in any maximal parabolic coset $gP$.
\end{defi}

The following set of inequalities is called weak relative Ancona inequalities and will be helpful below. They extend similar inequalities proved for hyperbolic groups by Ancona~\cite{Ancona} for $r=1$, by Gou\"ezel--Lalley~\cite{GouezelLalley} on co-compact Fuchsian groups for $r\leq R$ and by Gou\"ezel~\cite{Gouezel} in full generality for $r\leq R$. The version for relatively hyperbolic groups that we use here were first proved for $r=1$ by Gekhtman--Gerasimov--Potyagailo--Yang~\cite{GGPY} and then by Dussaule--Gekhtman~\cite{DG21} for $r\leq R$.

\begin{lemm}[{\cite[Theorem~1.1]{GGPY}, \cite[Theorem~1.6]{DG21}}]
\label{weakAncona}
Let $\Gamma$ be a relatively hyperbolic group and let $\mu$ be a finitely supported admissible and symmetric probability measure on $\Gamma$. Then, for every $c,\eta,L\geq 0$ there exists $C>0$ such that for every $r\leq R$, the following holds. For every $x,y,z\in \Gamma$, if $y$ has a distance at most $c$ to an $(\eta,L)$-transition point on a geodesic from $x$ to $z$, then
\[
\frac{1}{C}G(x,y|r)G(y,z|r)\leq G(x,z|r)\leq C G(x,y|r)G(y,z|r).
\]
\end{lemm}

Note that the constant $C$ is independent of $r\in [1,R]$. In other words, the Green function is roughly multiplicative along transition points on a geodesic. In both~\cite{DG21,GGPY}, these inequalities are formulated in terms of the Floyd distance, which is an appropriate rescaling of the word distance. However, the statement for transition points directly follows from~\cite[Corollary~5.10]{GePoGGD} which relates transition points with the Floyd distance. We also refer to~\cite[Section~9]{GGPY} for more details.

We will also use the following at some point.

\begin{lemm}[{\cite[Lemma~2.14]{YangPS}}]\label{ConCharLem}
There exist universal constants $\eta, L$ with the following property. Let $\gamma$ be a geodesic ray ending at a conical point $\xi\in\pG$. Then $\gamma$ contains a unbounded sequence of $(\eta,L)$-transition points $x_n$.
\end{lemm}


In the remainder of this section, we consider a finitely generated relatively hyperbolic group $\Gamma$. When speaking of a transition point, we mean an $(\eta,L)$-transition point with $(\eta,L)$ satisfying Lemma~\ref{ConCharLem}. We also fix a finitely supported symmetric and admissible probability measure $\mu$ on $\Gamma$.

\subsection{Busemann cocycles}
Given $x,y,z \in \Gamma$, let $B_z(x, y) := d(x, z) -d(y,z)$ and $K_z(x,y|r)=\frac{G(x,z|r)}{G(y,z|r)}$. The function $B_z$ is called the Busemann function associated with the distance $d$ at $z$.

Following~\cite{BlachereBrofferio}, we define the $r$-Green distance by
\[
d_r(x,y) = - \log F_r(x,y)=-\log \frac{ G(x,y|r)}{G(e,e|r)}.
\]
Then $K_z(x,y|r)=\rme^{-[d_r(x,z)-d_r(y,z)]}$ is the exponential of the Busemann function for the $r$-Green distance. We also write $|x^{-1}y|_r=d_r(x,y)$, and $|x^{-1}y|=d(x,y)$.

Consider the distance $\fkd_r$ defined for $x,y\in \Gamma$ by
\[
\fkd_r(x,y):=\omega_\Gamma(r)\left|x^{-1}y\right| +\left|x^{-1}y\right|_r.
\]

\begin{lemm}\label{quasiisometricdistances}
If $1\leq r<R$, then the distance $\fkd_r$ is proper and quasi-isometric to the word distance.
\end{lemm}

\begin{proof}
The proof is standard, but we write a complete argument for sake of completeness. First we prove that for every $r<R$, there exist $C_1> 0$ and $\alpha_1>0$ such that for every $x\in \Gamma$
\begin{equation}\label{upperboundquasiisometry}
G(e,x|r)\leq C_1\rme^{-\alpha |x|}.
\end{equation}
Since $\mu$ is finitely supported, there exists $c> 0$ such that
\[
G(e,x|r)=\sum_{n\ge c|x|}r^np_n(e,x).
\]
Moreover, by~\eqref{decayheatkernel},
\[
p_n(e,x)\leq R^{-n}.
\]
Therefore,
\[
G(e,x|r)\leq \sum_{n\ge c|x|}\left (\frac{r}{R}\right)^n\leq C_1\left (\frac{r}{R}\right)^{c|x|}.
\]
This proves~\eqref{upperboundquasiisometry}.

Second, we prove that for every $r\geq 1$, there exists $C_2>0$ and $\alpha_2>0$ such that for every $x\in \Gamma$,
\begin{equation}\label{lowerboundquasiisometry}
G(e,x|r)\geq C_2 \rme^{-\alpha_2 |x|}.
\end{equation}
Indeed, since the support of $\mu$ generates $\Gamma$, there exists a path $x_0=e,x_1,\dots,x_n=x$ such that $n\asymp |x|$ and $\mu(x_k^{-1}x_{k+1})>0$. In particular, we find that
\[
G(e,x|r)\geq G(e,x|1)\geq \mu\left(x_0^{-1}x_1\right)\dots \mu\left(x_{n-1}^{-1}x_n\right),
\]
which proves~\eqref{lowerboundquasiisometry}. We conclude that for $1\leq r<R$, the Green distance is quasi-isometric to the word distance and that $G(e,x|r)$ vanishes at infinity. Consequently, the distance $\fkd_r$ also is quasi-isometric to the word distance and satisfies that as $x$ goes to infinity, $\fkd_r(e,x)$ tends to infinity. In particular, any ball for $\fkd_r$ is contained in a larger ball for the word distance and thus is finite, so $\fkd_r$ is proper.
\end{proof}

\begin{rema}
According to~\cite[Lemma~2.1]{GouezelLalley}, for any non-amenable group $\Gamma$ and any finitely supported symmetric admissible probability measure $\mu$, $G(e,x|R)$ converges to 0 as $|x|$ goes to infinity. As a consequence, the distance $\fkd_R$ is proper, although it might not be quasi-isometric to the word distance.
\end{rema}

Define the corresponding Busemann cocycle
\begin{equation}\label{defBusemann}
\fkB_\xi(x, y;r) = \omega_\Gamma(r)B_\xi(x, y)-\log K_\xi(x,y|r)
\end{equation}



\begin{lemm}\label{BusemanEstLem}
There exists a constant $C>0$ with the following property. Let $\xi \in \pG$ be a conical point, and $x, y \in \Gamma$. There exists a neighborhood $V=V(x,y)$ of $\xi$ in $\Gamma\cup\pG$ such that for any $ z \in V \cap \Gamma$:
\begin{align*}
\left\lvert B_\xi(x, y)-B_z(x, y)\right| &\le C, \\
\left\lvert \log K_{\xi}(x, y|r)-\log K_{z}(x, y|r)\right| &\le C.
\end{align*}
\end{lemm}

\begin{proof}
The statement for $B_\xi$ is proved in~\cite[Lemma~2.20]{YangPS}. Also, by~\cite[Proposition~4.1]{DG21}, the Martin kernel $K_z(\cdot,e|r)=G(\cdot, z|r)/G(e, z|r)$ extends continuously to $K_\xi(\cdot,e|r)$ as $z$ converges to a conical limit point $\xi$. This follows from weak relative Ancona inequalities. In particular, $K_z(x,y|r)$ converges to $K_\xi(x,y|r)$ as $z$ converges to $\xi$, so the statement for $K_\xi$ also holds.
\end{proof}

As a consequence, the Busemann cocycle $\fkB_z(x,y)$ extends to a coarse cocycle $\fkB_\xi(x,y)$ at a conical point $\xi$. That is,
\[
\fkB_\xi(x,y) := \limsup_{z\to \xi} \fkB_\xi(x,y)
\]
does not depend on the choice of $z\to \xi$, up to a bounded additive error $C$ independent of $\xi$.

\subsection{Quasi-conformal densities}

A Borel measure $\mu$ on a topological Hausdorff space $T$ is \emph{regular} if $\mu(A) = \inf\{\mu(U): A \subset U, U$ is open$\}$ for any Borel set $A$ in $T$. It is called \emph{tight} if $\mu(A) =
\sup\{\mu(K): K \subset A, K$ is compact$\}$ for any Borel set $A$ in $T$. A finite Borel measure is called \emph{Radon} if it is tight and regular. It is well-known that all finite Borel measures on compact metric spaces are Radon, see~\cite[Theorem~1.1, Theorem~1.3]{Bill99}.

Denote by $\clM(\pG)$ the set of finite positive Radon measures on $\pG$. Then $\Gamma$ acts on $\clM(\pG)$ via $g_*\nu(A) = \nu(g^{-1}A)$ for any Borel set $A$ in $\pG$.

\begin{defi}
We call a map $x \mapsto \nu_x$ equivariant if for every $x,g\in \Gamma$, we have
\[
\nu_{gx}(A) = \nu_{x}\left(g^{-1}A\right)
\]
for every Borel set $A \subset\pG$.
\end{defi}



\begin{defi}
Let $\omega \in [0, \infty\mathclose{[}$. We call a $\Gamma$-equivariant map
\[
\nu: \Gamma \to
\clM(\pG),\quad x \mapsto \nu_x
\]
an
\emph{$\omega$-dimensional quasi-conformal density} if for any $x, y \in \Gamma$ the following holds
\begin{equation}\label{cdensitydef}
\frac{d\nu_{x}}{d\nu_{y}}(\xi) \asymp \rme^{-\omega B_\xi (x, y)}K_{\xi}(x,y|r),
\end{equation}
for $\nu_y$-a.e. points $\xi \in \pG$, where the implicit constant does neither depend on $x, y$, nor on $\xi$.
\end{defi}



We prove the following result.
\begin{lemm}
Let $\{\nu_x\}_{x\in\Gamma}$ be a $\sigma$-dimensional quasi-conformal density on $\pG$. Then the support of any $\nu_x$ is $\pG$.
\end{lemm}

\begin{proof}
By definition, the support $\supp(\nu_x)$ is a maximal closed subset such that any point in $\pG \setminus \supp(\nu_x)$ has an open neighborhood which is $\nu_x$-null. It is well-known that $\pG$ is a minimal $\Gamma$-invariant closed set, see for instance~\cite[Section~2]{Bo99}. Thus, it suffices to prove that the support of $\nu_x$ is $\Gamma$-invariant. This follows from equivariance and quasi-conformality, since $\nu_x(A)=\nu_{gx}(gA)$ for any Borel subset $A\subset \pG$ and $\nu_x$ and $\nu_{gx}$ are absolutely continuous with respect to each other.
\end{proof}



As explained in the introduction, we associate the following Poincaré series to $\mu$ and to the word distance, by setting
\begin{align*}
\Theta_\Gamma({r, s}) := \sum_{x\in \Gamma} G(e, x |r) \rme^{-s|x|} = \sum_{n\ge 0} H_r(n) \rme^{-sn} =G(e,e|r)\cdot \sum_{x\in \Gamma} \rme^{-s|x|-|x|_r}
\end{align*}
where we recall that $H_r(n)=\sum_{x\in S_n} G(e,x|r)$ and that the
\emph{critical exponent} $\omega_\Gamma(r)$ is defined by
\[
\omega_\Gamma(r) = \limsup\limits_{n \to \infty} \frac{1}{n} \log H_r(n).
\]
The group $\Gamma$ is of \emph{divergent} (resp.
\emph{convergent}) type {for the Green function} if $\Theta_\Gamma({r,s})$ is divergent (resp. convergent) at $s=\omega_\Gamma(r)$.

Recall $\fkd_r(x,y):=\omega_\Gamma(r)|x^{-1}y| +|x^{-1}y|_r$.
\begin{lemm}\label{NewPoincareSeries}
The series $\clP_\Gamma$ defined as
\[
\forall s>0,\;\clP_{\Gamma}(s):=\sum\limits_{x \in \Gamma} \rme^{- s \fkd_r(x,y)}
\]
has critical exponent $1$, and the divergence of $\clP_{\Gamma}(s)$ at $s=1$ is equivalent to that of $\Theta_\Gamma({r,s})$ at $s=\omega_\Gamma(r)$.
\end{lemm}

\begin{proof}
By~\cite[Lemma~2.1]{GouezelLalley}, $G(e,x|R)$ goes to 0 as $|x|$ tends to infinity. Thus, $G(e,x|r)\leq G(e,x|R)\leq 1$ for large enough $x$. In particular, we see that for $s> 1$,
\[
\sum_{x\in \Gamma}G(e,x|r)^s\rme^{-s\omega_\Gamma(r) |x|}\leq C \sum_{x\in \Gamma}G(e,x|r)\rme^{-s\omega_\Gamma(r) |x|}.
\]
The right-hand side converges, so the left-hand side also converges. Conversely, for $s<1$,
\[
\sum_{x\in \Gamma}G(e,x|r)\rme^{-s\omega_\Gamma(r) |x|}\leq C \sum_{x\in \Gamma}G(e,x|r)^s\rme^{-s\omega_\Gamma(r) |x|}.
\]
The left-hand side diverges, so the right-hand side also diverges. Thus, the critical exponent of $\clP_\Gamma(s)$ is 1.

Also, note that $\clP_\Gamma(1)=\Theta_\Gamma(r,\omega_\Gamma(r))$, so the second conclusion of the lemma follows.
\end{proof}

Write $\mu(f) = \int f d\mu$ for a continuous function $f \in C(\pG)$. We endow $\clM(\pG)$ with the weak topology. That is, a sequence $\mu_n \in \clM(\pG)$ converges to $\mu$ if $\mu_n(f) $ converges to $ \mu(f)$ for any $f \in C(\pG)$. Equivalently, by the Portmanteau Theorem~\cite[Theorem~2.1]{Bill99}, $\mu_n$ converges to $\mu$ if $\liminf_{n \to \infty}\mu_n(U) \ge \mu(U)$ for any open set $U \subset \pG$.

We start by constructing a family of measures $\{\nu_x^s\}_{x\in \Gamma}$ supported on $\Gamma$ for any $s >1$. First, assume that $\clP_\Gamma(s)$ is divergent at $s=1$. Set
\[
\nu^s_x = \frac{1}{\clP_{\Gamma}(s)} \sum\limits_{z \in \Gamma} \rme^{-s\fkd_r(x,z)} \cdot \dirac{z},
\]
where $s >1$ and $x \in \Gamma$. Note that $\nu^s_x$ is a probability measure.


On the contrary, assume that $\clP_{\Gamma}(s)$ is convergent at $s=1$, Patterson introduced in~\cite[Lemma~3.1]{Patt} a monotonically increasing function $H: \R_{\ge 0} \to \R_{\ge 0}$ with the following property:
\begin{equation}\label{patterson}
\forall \epsilon >0, \;\exists t_\epsilon,\quad \forall t > t_\epsilon,\quad \forall a>0,\quad  H(a+t) \le \exp(a\epsilon) H(t).
\end{equation}
and such that the following modified series
\[
\clP_{\Gamma}'(s):=\sum\limits_{x \in \Gamma} H(\fkd_r(x,z)) \cdot \rme^{-s\fkd_r(x,z)}
\]
is divergent for $s \le 1$ and convergent for $s>1$.
Then define measures as follows:
\[
\nu^s_x = \frac{1}{\clP_{\Gamma}'(s)} \sum\limits_{z \in \Gamma} \rme^{-s \fkd_r(x,z)} \cdot H(\fkd_r(x,z)) \cdot \dirac{z},
\]
where $s >1$ and $x \in \Gamma$.

Choose $s_i \searrow 1$ such that $\nu_x^{s_i}$ are convergent in $\clM(\pG\cup\Gamma)$ for all $x\in \Gamma$. Let $\nu_x = \lim \nu_x^{s_i}$ be the limit measures, which are called \emph{Patterson--Sullivan measures} associated with the Poincaré series $\clP_\Gamma$. Note that forcing the Poincaré series to be divergent at 1, we have $\nu_x(\pG) = 1$. In the sequel, we write PS-measures as shorthand for Patterson--Sullivan measures. We also write $\partial\Gamma^{\con}$ for the set of conical limit points in the Bowditch boundary.



\begin{prop}\label{prequasiconf}
The PS-measures $\{\nu_x\}_{x\in \Gamma}$ on the Bowditch boundary are absolutely continuous with respect to each other and satisfy
\begin{align}
\forall \xi\in \pG ,\quad  \frac{d\nu_x}{d\nu_y}(\xi) &\ge \rme^{-\fkd_r(x,y)},\label{LowerBoundEQ}\\
\forall \nu_y\; \text{a.e.}\; \xi\in \pG^{\con},\quad  \frac{d\nu_{x}}{d\nu_{y}}(\xi) &\asymp \rme^{-\omega_\Gamma(r) B_\xi (x, y)}K_{\xi}(x,y|r),\label{cdensity}
\end{align}
where the implicit constant is independent of $x,y$ and $\xi$.
\end{prop}


\begin{rema}
If $\Gamma$ is of divergent type for Green function, then Theorem~\ref{divergencetypeandatoms} below says that PS-measures have no atoms on Bowditch boundary and give full measure to conical limit points, so~\eqref{cdensity} holds for $\nu_y$-a.e. $\xi\in \pG$. In this case, $\nu$ is an $\omega_\Gamma(r)$-dimensional quasi-conformal density.
\end{rema}
\begin{proof}
Since $\fkd_r$ satisfies the triangle inequality and $\lim_{t\to\infty} \frac{H(a+t)}{H(t)}=1$, we see that $\{\nu_x: x\in \Gamma\}$ are absolutely continuous with respect to each other,
\begin{align*}
\rme^{-\fkd_r(x,y)}\le \frac{d\nu_x}{d\nu_y}(\xi) \le \rme^{\fkd_r(x,y)}.
\end{align*}


We now verify the quasi-conformality at conical limit points. We only consider the case where $\Gamma$ is of convergent type, the divergent type being simpler. Let $\epsilon> 0$ and $t_\epsilon$ the number be given by~\eqref{patterson} for the function $H$. Let $\phi=\frac{d\nu_x}{d\nu_y}$ be the Radon--Nikodym derivative, uniquely defined up to a $\nu_y$-null set. Let $\xi \in \pG$ be a conical limit point and consider the open neighborhood $V$ of $\xi$ and the uniform constant $C$ given by Lemma~\ref{BusemanEstLem}.

Let $f$ be a continuous function supported in $V$. One can choose $V$ also such that $\fkd_r(y, z) > t_\epsilon$ for all $z \in V$. If $z \in V$ satisfies $\fkd_r(x, z) >\fkd_r(y, z)$, then we have
\begin{align*}
H(\fkd_r(x,z)) &= H(\fkd_r(x,z)-\fkd_r(y,z)+\fkd_r(y,z)) \\
& \le \rme^{\epsilon[\fkd_r(x,z)-\fkd_r(y,z)]} \cdot H(\fkd_r(y,z)) \\
& \le \rme^{\epsilon (C + \fkB_\xi(x, y))}\cdot H(\fkd_r(y, z)).
\end{align*}
Since $H$ is increasing, we have
\begin{equation}\label{prequasiconfeq}
C_\epsilon^{-1} H(\fkd_r(y, z))\le H(\fkd_r(x, z)) \le C_\epsilon H(\fkd_r(y, z)),
\end{equation}
where $C_\epsilon= \rme^{\epsilon (C + \fkB_\xi(x, y))}>1$ depends on $\epsilon, C$ and $ (x,y)$, but not on $z\in V$. Note that $C_\epsilon\to 1$ as $\epsilon\to 0$. By symmetry, the conclusion~\eqref{prequasiconfeq} also holds if $\fkd_r(x, z) < \fkd_r(y, z)$ for $z \in V$.

Using~\eqref{prequasiconfeq}, we get the following estimates. First,
\begin{align*}
\nu^s_x(f) & = \frac{1}{\clP_{\Gamma}(s)} \sum\limits_{z \in V}
\rme^{-s \fkd_r(x,z)} H(\fkd_r(x,z)) {f(z)} \\
& \leq C_\epsilon \rme^{-s\fkB_{z}(x, y))}\cdot \frac{1}{\clP_{\Gamma}(s)}
\sum\limits_{z \in V}
\rme^{-s\fkd_r(y, z)} H(\fkd_r(y, z)) {f(z)}\\
& \asymp C_\epsilon \rme^{-s\fkB_\xi(x, y)} \nu^s_y(f)
\end{align*}
and second
\begin{align*}
\nu^s_x(f) & = \frac{1}{\clP_{\Gamma}(s)} \sum\limits_{z \in V}
\rme^{-s \fkd_r(x,z)} H(\fkd_r(x,z)) {f(z)} \\
& \geq C_\epsilon^{-1} \rme^{-s\fkB_{z}(x, y))}\cdot \frac{1}{\clP_{\Gamma}(s)}
\sum\limits_{z \in V}
\rme^{-s\fkd_r(y, z)} H(\fkd_r(y, z)) {f(z)}\\
& \asymp C_\epsilon \rme^{-s\fkB_\xi(x, y)} \nu^s_y(f)
\end{align*}
where the implicit constants depend only on $C$ but not on $\epsilon$. Letting $s\to 1$, we get
\[
C_\epsilon^{-1} \rme^{-\fkB_\xi(x, y)}\nu_y(f)\prec \nu_x(f) \prec C_\epsilon \rme^{-\fkB_\xi(x, y)} \nu_y(f)
\]
for any continuous function $f$ supported in $V$. As $\epsilon\to 0$, $C_\epsilon\to 1$, hence it follows that
\[
\phi(\xi) \asymp \rme^{-
\fkB_\xi (x, y)}
\]
for $\nu_y$-a.e. conical limit point $\xi \in \pG$.
\end{proof}

\subsection{Shadow Lemma}

The key observation in the theory of PS-measures is the Sullivan Shadow lemma shedding some light on the relation between the PS-measure and the geometric properties of the boundary.

Recall that $\partial \Gamma$ is a visual boundary for the Cayley graph $\Cay(G,S)$: any two distinct points $x,y\in \Cay(G,S)\cup \partial \Gamma$ can be connected by a geodesic. See~\cite[Proposition~2.4]{GePoJEMS}.

\begin{defi}
Let $C>0$ and $x \in\Gamma$. The \emph{shadow $\Pi_C(x)$} at $x$ is the set of points $\xi \in \pG$ such that there exists some geodesic $\gamma=[e, \xi]$ intersecting $B(x, C)$.

The \emph{partial shadow $ \Psi_{C}(x)$} at $x$ is the set of points $\xi \in \pG$ such that some geodesic $[e,\xi]$ contains a transition point $C$-close to $x$.
\end{defi}

We now state the Shadow lemma in our context, whose proof follows closely the proofs of~\cite[Lemmas~4.1 \& 4.2]{YangPS} with Lemma~\ref{BusemanEstLem} replacing Lemma~2.19 there. Let us denote by $\Psi^{\con}_C(g)$ the set of all conical limit points in $\Psi_C(g)$.

\begin{lemm}[Shadow lemma]\label{ShadowLem}
Let $\{\nu_x\}_{x \in \Gamma}$ be an $\omega_\Gamma(r)$-dimensional PS measures on the Bowditch boundary $\pG$. Then there exists $C_0> 0$ such that for any $C\ge C_0$ and $x\in \Gamma$ the following two inequalities hold
\begin{align}
\rme^{-\omega_\Gamma(r) |x|}G(e,x|r) &\prec_C \, \nu_e(\Psi_{C}(x))\le \nu_e(\Pi_{C}(x)),\label{LowBoundEQ}\\
\nu_e(\Psi^{\con}_{C}(x)) &\prec_C \, \rme^{-\omega_\Gamma(r) |x|} G(e,x|r).\label{UpperBoundEQ}
\end{align}
\end{lemm}

\begin{rema}
If $\nu_e$ has no atoms at parabolic points which form a countable subset of the Bowditch boundary, then we obtain the full strength of the partial shadow lemma without having to restrict our attention to conical points. The upper bound~\eqref{UpperBoundEQ} for the partial shadow uses the relative Ancona inequalities (Lemma~\ref{weakAncona}), while it is unknown whether the upper bound holds for the usual shadow.
\end{rema}

\begin{proof}
Let $F$ be a set of three loxodromic elements with pairwise disjoint fixed points. For each $f\in F$, let $\alpha:=\bigcup_{i\in \Z} f^i[e,f]$ be an $\langle f\rangle$-invariant quasi-geodesic between two fixed points $f_-,f_+\in\pG$. Let $U_f\subset \pG$ be an open neighborhood of $f_+$ so that for any $\eta\in U_f$, the projection of $\eta$ to the axis $\alpha$ has a distance to $e$ at least $C$. By~\cite[Lemma~2.4]{DWY}, for any $x\in \Gamma$, there exists $f\in F$ so that $[x^{-1},\eta]$ contains a transition point $C$-close to $e$. Thus, $U_f\subset x^{-1}\Psi_{C}(x)$.

As $\Gamma$ acts minimally with a dense orbit in $\pG$, the $\Gamma$-orbit of any open set $U\subset \pG$ covers $\pG$, so $U$ have positive $\nu_e$-measure. Hence, setting
\[
D=\min\{\nu_e(U_f): f\in F\}>0
\]
which is independent of $x$, we have
\[
\nu_e\left(x^{-1}\Psi_{C}(x)\right) \ge D.
\]
Since $\nu_x $ is equivariant, the lower bound in~\eqref{LowerBoundEQ} implies
\begin{align*}
\nu_e(\Psi_{C}(x))={\nu_{x^{-1}}\left(x^{-1}\Psi_{C}(x)\right)}&\ge \rme^{-\omega_\Gamma(r) |x|}\frac{G(e,x|r) }{G(e,e|r)}\cdot {\nu_e\left(x^{-1}\Psi_C(x)\right)}\\
&\ge D \rme^{-\omega_\Gamma(r) |x|}\frac{G(e,x|r)}{G(e,e|r)}.
\end{align*}
Since $G(e,e|r)$ is bounded by $G(e,e|R)$, this concludes the proof of the lower bound.

For any $\xi \in x^{-1} \Psi^{\con}_C(x)$, there is a geodesic $\gamma$ from $x^{-1}$ to $\xi$ which intersects $B(e, C)$ and contains a transition point. Thus, $|B_\xi(x^{-1},e)
- d(x^{-1}, e)| \le 2C.$ By the relative Ancona inequalities (Lemma~\ref{weakAncona}), there is a constant $C_1$ independent of $r$ such that
\[
K_\xi\left(x^{-1}, e|r\right) = \lim_{z\to \xi}\frac{G\left(x^{-1},z|r\right)}{G(e,z|r)}
\le C_1 G(e,x|r).
\]
Also, by~\eqref{cdensity} there is a constant $C_2>0$ such that for $\nu_e$-a.e. conical limit point $\xi \in \pG^{\con}$,
\[
\frac{d\nu_{x^{-1}}}{d\nu_e}(\xi) \le C_2 \rme^{-\omega_\Gamma(r) B_\xi\left(x^{-1},e\right)} K_\xi\left(x^{-1}, e| r\right).
\]
Combining together the above estimates, we have
\begin{align*}
\nu_e(\Psi^{\con}_C(x)) & = \nu_{x^{-1}}\left(x^{-1}\Psi^{\con}_C(x)\right) \\
& \le C_2\int_{x^{-1}\Psi^{\con}_C(x)} \rme^{-\omega_\Gamma(r) B_\xi\left(x^{-1},e\right)} K_\xi\left(x^{-1}, e|r\right) d \nu_e(\xi) \\
& \leq C_1C_2 \rme^{2C\omega_\Gamma(r)} \cdot \rme^{-\omega_\Gamma(r) |x|} G(e,x|r),
\end{align*}
which finishes the proof of the upper bound.
\end{proof}



\begin{prop}\label{PosMeasonConcialProp}
Suppose that $\nu_e$ gives positive measure to the set of conical limit points. Then $\Gamma$ is of divergent type for the Green function.
\end{prop}
\begin{proof}
List $\Gamma = \{x_1, \dotso, x_i, \dots\}$ such that for all $i$, $|x_i|\le | x_{i+1}|$. Let $C_0$ be given by Lemma~\ref{ShadowLem}. For any $C>C_0$, set
\begin{equation}\label{conicalset}
A_{C} := \bigcap_{n=1}^{\infty} \bigcup_{i=n}^{\infty}
\Psi_{C}^{\con}(x_i).
\end{equation}
By Lemma~\ref{ConCharLem}, we have $\partial^{\con}\Gamma =A_{C}$. In other words, any conical limit point is shadowed infinitely many times by elements of $\Gamma$.

We claim that $\Theta_{\Gamma}(r,\omega_\Gamma(r))$ is divergent. Recall that
\[
\Theta_{\Gamma}(r,\omega_\Gamma(r)) = \sum\limits_{x \in \Gamma} \rme^{-\omega_\Gamma(r)|x|}G(e,x|r).
\]
By Lemma~\ref{ShadowLem}, we see that
\begin{align*}
\sum_{|x|\ge n} \rme^{-\omega_\Gamma(r)|x|}G(e,x|r) & \succ \sum_{|x|\ge n} \nu_e(\Psi^{\con}_{C}(x)) \\
&\succ \nu_e\left(\bigcup_{|x|\ge n} \Psi^{\con}_C(x)\right) \succ \nu_e(A_{C}) > 0.
\end{align*}
This is true for all $n > 0$ and $\nu_e(A_{C})$ is independent of $n$. Thus, $\Theta_{\Gamma}(r,s)$ is indeed divergent at $s=\omega_\Gamma(r)$.
\end{proof}


\begin{theo}\label{divergencetypeandatoms}
If $\Gamma$ is of divergent type for Green function, then $\nu_e$ has no atoms. Otherwise, $\nu_e$ is purely atomic and supported on the set of parabolic points.
\end{theo}
\begin{proof}
First of all, the otherwise statement follows from Proposition~\ref{PosMeasonConcialProp}. Indeed, if $\Gamma$ is of convergent type, then $\nu_e$ cannot give positive measure to conical limit points, hence it is supported on the set of parabolic limit points. Since this set is countable, $\nu_e$ is necessarily purely atomic.

Assume now that $\Gamma$ is of divergent type. Let $q\in \pG$ be a bounded parabolic point so that the stabilizer $P\in \bbP$ acts co-compactly on $\pG\cup \Gamma \setminus \{q\}$. Let $K\subset \Gamma\cup\pG\setminus \{q\}$ be a compact fundamental domain.
For a point $y\in \Gamma$, we let $\pi_P(y)$ be the set of nearest point-projections
\[
\pi_P(y)=\{p\in P: d(y,p)=d(y,P)\}.
\]
Define $\pi_P(A):=\bigcup_{a\in A} \pi_P(a)$. As $\partial K$ is disjoint with $\partial P=\{q\}$, the shortest projection $Z:=\pi_P(K\cap \Gamma)$ has bounded diameter by~\cite[Proposition~3.3]{GP16}. By enlarging~$Z$, assume without loss of generality that $1\in Z$.

Note that a maximal parabolic $P$ has the contracting property by~\cite[Proposition~8.5]{GP16}: any geodesic $[x,y]$ with large projection to $P$ is uniformly close to $\pi_P(x)$ and $\pi_P(y)$. Thus, for any $y\in \Gamma\cap K$ and $p\in P$, any geodesic $[e,py]$ passes within uniformly bounded distance of $p\in pZ$. Consequently, $|p y| \simeq |p| + |y|$. Moreover, $[e,py]$ exits $P$ with bounded distance to $pZ$, so $p$ is within a bounded distance of a transition point on $[e,py]$. Now, by the relative Ancona inequalities in Lemma~\ref{weakAncona}, $G(e, p y|r) \asymp G(e, p|r) G(e, y|r)$. Here the implicit constants in $\simeq$ and $\asymp$ are independent of $y$ and $p$.



We can estimate the $\nu_e^s$-measure of an open neighborhood
\[
U_n= \{q\}\cup \{pK: p\in P; |p|\ge n\}
\]
as follows
\[
\nu_e^s(U_n) \le \sum_{|p|\,>\,n} \nu_e^s(pK) \le {\frac{1}{\clP_{\Gamma}(s)}} \sum_{|p|\,>\,n}\; \sum_{y\in K} \rme^{-s\omega_\Gamma(r)|p y|} [G(e, p y|r)]^s,
\]
hence
\begin{align*}
\nu_e^s(U_n) &\prec {\frac{1}{\clP_{\Gamma}(s)}} \sum_{|p|\,>\,n}\;\sum_{y\in K} \rme^{-s\omega_\Gamma(r)|p|} \rme^{-s\omega_\Gamma(r)|y|} [G(e,y|r)]^s [G(e,p|r)]^s\\
&\prec \nu_e^s(K) \sum_{|p|\,>\,n} \rme^{-s\omega_\Gamma(r) |p|} [G(e,p|r)]^s.
\end{align*}


By~\cite[Corollary~3.9]{DWY}, $\Theta_P(r,\omega_\Gamma(r))$ is convergent. Also, by the Portmanteau theorem~\cite[Theorem~2.1]{Bill99}, $\limsup_{s\to 1}\nu_e^s(K)\leq \nu_e(K)$. Letting $s\to 1$ and then $n\to\infty$, we see that $\nu_e(U_n)\to 0$. Thus, $\nu_e$ has no atoms at parabolic limit points. By~Lemma~\ref{ConCharLem}, a conical limit point $\xi$ is contained in infinitely many partial shadows $\Psi_{C}(x_n)$. By Lemma~\ref{ShadowLem}, as $x_n\to \xi$, the $\nu_e$-measure of $\Psi_C(x_n)$ tends to $0$, so conical limit points are not atoms as well.
\end{proof}


\section{Convergent Poincaré series and applications}
\label{SectionRHG}

This final section is devoted to the proof of Theorem~\ref{mainthm} and Corollary~\ref{transitionphaseCor}. One important notion that was coined in~\cite{DWY} is the \emph{parabolic gap for the Green function} whose definition we now recall. Let $\mu$ be a probability measure on $\Gamma$, let $\bbP_0$ be a finite set of representatives of conjugacy classes of maximal parabolic subgroups and let $P\in \bbP_0$. We set
\[
H_{P,r}(n)=\sum_{x\in S_n\cap P}G(e,x|r)
\]
and
\[
\omega_P(r)=\limsup_{n\to \infty}\frac{1}{n}\log H_{P,r}(n).
\]
\begin{defi}
If $\omega_P(r)<\omega_\Gamma(r)$, we say that $\Gamma$ has a \emph{parabolic gap} along $P$ for the Green function at $r$. If for every $P$, for every $r\in (1,R]$, $\omega_P(r)<\omega_\Gamma(r)$, then we say that $\Gamma$ has a \emph{parabolic gap} for the Green function.
\end{defi}


One of the consequences of having a parabolic gap for the Green function is that $H_r(n)$ is roughly multiplicative and the Green function has purely exponential growth. Namely, by~\cite[Theorem~1.8]{DWY}, if $\Gamma$ is a non-elementary relatively hyperbolic group and if $\mu$ is a finitely supported admissible and symmetric probability measure on $\Gamma$ such that the parabolic condition for the Green function holds, then for every $1<r\leq R$, there exist $C$ and $C'$ such that for all $n$,
\begin{equation}\label{roughmultiplicativity}
\frac{1}{C}H_{r} (n+m)\leq H_r(n)H_r(m)\le C H_{r} (n+m)
\end{equation}
and
\begin{equation}\label{purelyexponentialgrowth}
\frac{1}{C'}\rme^{n\omega_\Gamma(r)}\leq H_r(n)\leq C' \rme^{n\omega_\Gamma(r)}.
\end{equation}
As proved in~\cite{SWX}, these two properties~\eqref{roughmultiplicativity} and~\eqref{purelyexponentialgrowth} hold for all hyperbolic groups. It was proved in~\cite{DWY} that if maximal parabolic subgroups are amenable and if $r<R$, then the parabolic gap condition holds, hence so do the properties~\eqref{roughmultiplicativity} and~\eqref{purelyexponentialgrowth}. Under additional assumptions on the random walk, this was also proved at~$R$. Among the unanswered problems in~\cite{DWY} are the following questions. Does the parabolic gap condition holds at~$R$ as soon as maximal parabolic subgroups are amenable? Does there exist examples of relatively hyperbolic groups endowed with a finitely supported admissible symmetric probability measures such that the properties~\eqref{roughmultiplicativity} and~\eqref{purelyexponentialgrowth} fail? As particular cases of our work, we answer these two questions here.

\subsection{First return to maximal parabolic subgroups}

We gather the results of Section~\ref{Sectiongrowthrateat1} to prove that whenever maximal parabolic subgroups are amenable, the parabolic gap for the Green functions holds.


Consider the first return kernel $p_{r,P}$ to $P$ for $r \mu$ defined by
\[
p_{r,P}(x,y)=\sum_{n\ge 1}\;\sum_{z_1, \ldots,z_{n-1}\notin P} r^n p~(x,z_1)p(z_1,z_2) \dots p(z_{n-1},y).
\]
Denote by $p_{r, P}^{(n)}$ the $n^{\rm th}$ convolution power of this transition kernel and by $G_{r,P}$ the associated Green function. It turns out that for every $x$, $y\in P$, we have $G_{r,P}(x,y|1)=G(x,y|r)$; see~\cite[Lemma~4.4]{DG21}. Set $t_{r, P} = \sum_{x \in P} p_{r, P}(e, x)$. Then $t_{r, P}^{-1} p_{r, P}$ is a symmetric admissible and $P$-invariant transition kernel, thus it defines a random walk on $P$.

\begin{prop}\label{omegaPlesst}
Let $\Gamma$ be a relatively hyperbolic group and let $P\in \bbP_0$. Consider an admissible and symmetric probability measure $\mu$ on $\Gamma$. If $t_{r, P} \leq 1$, then we have $\omega_P(r) \leq 0$; in particular, $\omega_P(r) < \omega_{\Gamma}(r)$.
\end{prop}

\begin{proof}
Denote $t = t_{r, P}$ for simplicity and let $G_t$ be the Green function associated with $t^{-1} p_{r, P}$. By Proposition~\ref{upperboundgrowthrateat1}, $\sum_{x \in P, \, |x| = n} G_t(e, x | 1)$ has growth rate at most 0. Since $t \leq 1$ we have for every $x\in P$ that
\[
G(e, x | r) = G_{r, P}(e, x | 1) = G_t(e, x | t) \leq G_t(e, x | 1).
\]
Therefore $\omega_P(r) \leq 0$.
\end{proof}

\begin{prop}\label{propamenableparabolic}
Let $\Gamma$ be a relatively hyperbolic group and let $P$ be a maximal parabolic subgroup. Consider an admissible and symmetric probability measure $\mu$ on $\Gamma$. If $P$ is amenable, then for every $r\leq R$, $\omega_P(r)\leq 0$.
\end{prop}



\begin{proof}
Since $G(e, x | r) < \infty$, we deduce that the spectral radius of $p_{r, P}$ is at most $1$. By~\cite[Corollary~12.5]{Woessbook} and the fact that $P$ is amenable, the spectral radius of $t_{r, P}^{-1} p_{r, P}$ is $1$ and hence $t_{r, P} \leq 1$. The result follows from Proposition~\ref{omegaPlesst}.
\end{proof}

Note that we do not need to assume that $\mu$ is finitely supported in this proposition, although we need this assumption in the following corollary, which also relies on~\cite[Theorem~1.8]{DWY} mentioned above, where the assumption is crucially used.

\begin{coro}\label{coramenableparabolic}
Let $\Gamma$ be a relatively hyperbolic group endowed with a finitely supported admissible and symmetric probability measure $\mu$. Assume that maximal parabolic subgroups of $\Gamma$ are amenable. Then $\Gamma$ has a parabolic gap for the Green function and so~\eqref{roughmultiplicativity} and~\eqref{purelyexponentialgrowth} hold.
\end{coro}



Note that $t_{1, P}$ is the probability that the random walk associated to $\mu$ eventually returns to $P$. We see that $t_{1, P} < 1$ for all $P \in \bbP_0$. Otherwise, the random walk would visit $P$ infinitely many times with positive probability, which in turn would imply that it accumulates at the parabolic limit points fixed by $P$. This would contradict the fact that the random walk almost surely converges to a conical limit point~\cite[Theorem~9.14]{GGPY}.



More generally, consider a branching random walk on $\Gamma$ whose step distribution is given by $\mu$ and with mean offspring $r$. Following~\cite{CandelleroGilchMuller}, by collecting all particles returning to $P$, one gets a Galton--Watson process with mean offspring $t_{r,P}$, see precisely the proof of~\cite[Proposition~4.3]{CandelleroGilchMuller}. Consequently, the branching random walk returns to $P$ infinitely many times if and only if $t_{r,P}>1$. In such case, the branching random walk must accumulate in $P$. Moreover, the same arguments as in~\cite[Proposition~4.4]{CandelleroGilchMuller} show that there are almost surely infinitely many cosets $gP$, such that the branching random walk accumulates in $gP$. On the contrary, if $t_{r,P}\leq 1$, then for all coset $gP$, the branching random walk eventually leaves $gP$. This follows from the fact that $g$ is visited finitely many time almost surely and that starting a branching random walk at $g$, it comes back to $gP$ only finitely many times, see also~\cite[Proposition~4.5]{CandelleroGilchMuller} for a more detailed proof.

Define now
\[
r_{\clP} = \sup \{ r > 1 \colon t_{r, P} \leq 1 \quad\text{for all } P \in \bbP_0 \}.
\]
Then for $1<r\leq r_\clP$, we have $\omega_P(r)\leq 0$ for all $P\in \bbP_0$, hence $\Gamma$ has a parabolic gap for the Green function and $H_r(n)$ has purely exponential growth. Moreover, $r_\clP$~is the transition for the branching random walk spending infinite times in maximal parabolic subgroups: if $r \leq r_\clP$, then the branching random walk eventually leaves every coset $gP$, $P\in \bbP_0$, while if $r>r_\clP$, it accumulates in infinitely many cosets $gP$ for at least one of the $P\in \bbP_0$.

Recall that the limit set $\Lambda$ of a branching random walk inside the Bowditch boundary $\partial\Gamma$ is the set of accumulation points in $\partial\Gamma$ of the trace of $\BRW(\Gamma,\nu,\mu)$, which is the set of elements of $\Gamma$ that are ever visited by the branching random walk. We take the occasion to derive the following consequence which sheds some light on the geometry of the limit set.

\begin{prop}\label{propparaboliclimitset}
Let $\Gamma$ be a relatively hyperbolic group endowed with a finitely supported admissible and symmetric probability measure $\mu$. Consider a probability measure $\nu$ on $\bbN$ with mean $r\leq R$. Let $\Lambda$ be the limit set inside the Bowditch boundary of $\Gamma$ of the branching random walk associated with $\mu$ and $\nu$.
\begin{enumerate}
\item\label{prop5.5.1} If $r<r_\clP$, then almost surely, $\Lambda$ does not contain any parabolic limit point.
\item\label{prop5.5.2} If $r>r_\clP$, then almost surely, $\Lambda$ contains an infinite number of parabolic limit points.
\end{enumerate}
\end{prop}

Note that $\Lambda$ cannot contain all parabolic limit points in case~\eqref{prop5.5.2}. Otherwise, since parabolic limit points are dense, $\Lambda$ would coincide with the whole Bowditch boundary, but this is impossible since its Hausdorff dimension with respect to the shortcut distance is at most half the Hausdorff dimension of the whole boundary by~\cite[Theorem~1.1, Theorem~1.2]{DWY}.

This result is a consequence of the following one. For $x\in \Gamma$ and $C\geq0$, we denote by $\Omega(x,C)$ the \emph{partial cone} at $x$ of width $C$, which is the set of points $y\in \Gamma$ such that $x$ is within $C$ of a transition point on a geodesic from $e$ to $y$.

\begin{prop}\label{proppartialcone}
Let $\Gamma$ be a relatively hyperbolic group endowed with a finitely supported admissible and symmetric probability measure $\mu$. For every $r\in [1,R]$, there exist $\beta>0$ with the following property. Consider a probability measure $\nu$ on $\bbN$ with mean $r\leq R$. For any $x\in \Gamma$, the probability that the branching random walk visits $\Omega(x,C)$ is at most $C_1(1+|x|^\beta)G(e,x|r)$, where $C_1$ is a constant.
\end{prop}

The proof of this proposition relies mostly on material from~\cite{DWY} and its proof is postponed to the Appendix~\ref{app}.

\begin{proof}[{Proof of Proposition~\ref{propparaboliclimitset}}]
As we saw above, if $r>r_\clP$, then the branching random walk accumulates in infinitely many cosets $gP$ for at least one $P\in \bbP_0$. In particular, the limit set $\Lambda$ contains all parabolic limit points fixed by $gPg^{-1}$, for every such coset $gP$. Thus, we only need to prove~\eqref{prop5.5.1} to conclude and we assume that $r<r_\clP$.

Recall that a sequence $x_n$ in $\Gamma$ converges to a parabolic limit point $\xi$ fixed by $gPg^{-1}$, $P\in \bbP_0$, if and only if the sequence of projections of $x_n$ on $gP$ tends to infinity. Thus, $\xi\in \Lambda$ if and only if the branching random walk visits infinitely many $\Omega(x,C)$, with $x\in gP$. Fix $g\in \Gamma$ and $P\in \bbP_0$. Denote by $A_n$ the event
\[
A_n=\big\{\BRW(\Gamma,\nu,\mu)\text{ visits } \Omega(x,C) \text{ for some } x\in gP,\text{ with } \left|g^{-1}x\right|=n\big\}.
\]
Then, by Proposition~\ref{proppartialcone},
\[
\bfP(A_n)\leq C_1(1+n)^\beta \sum_{x\in P,|x|=n}G(e,gx|r).
\]
Since $g$ is a transition point on a relative geodesic from $e$ to $gx$, by relative Ancona inequalities we get
\[
\bfP(A_n)\leq C_2(1+n)^\beta H_{P,r}(n).
\]
Since $t_{r,P}<1$, $\omega_P(r)<0$ by Proposition~\ref{omegaPlesst} and so $H_{P,r}(n)$ decays exponentially fast as $n$ goes to infinity. This proves that
\[
\sum_{n}\bfP(A_n)<\infty.
\]
By the Borel--Cantelli lemma, we deduce that almost surely, the branching random walk only visits finitely many $\Omega(x,C)$, $x\in gP$, hence $\xi\notin \Lambda$. Since parabolic limit points are countable, this settles the proof.
\end{proof}



The critical case $r=r_\clP$ remains open. In the context of free products, the branching random walk needs to visit $gP$ infinitely many times in order to accumulate at $\xi$. As a consequence, the authors of~\cite{CandelleroGilchMuller} prove in Propositions~4.3,~4.4 and~4.5 that for $r=r_\clP$, almost surely, $\Lambda$ does not contain any parabolic limit point. However, for arbitrary relatively hyperbolic groups, the branching random walk can visit infinitely many partial cones $\Omega(x,C)$, $x\in gP$, without entering $gP$ at all. Thus, we might need new material to figure this critical case.



\subsection{Convergent Poincaré series}
\label{SectionconvergentPoincare}

Let $\Gamma$ be a relatively hyperbolic group and let $\mu$ be a finitely supported symmetric admissible probability measure on $\Gamma$. We consider the Poincaré series $\Theta_\Gamma(r,s)$ defined above and for $P$ a maximal parabolic subgroup, the Poincaré series $\Theta_P(r,s)$ defined~by
\[
\Theta_P(r,s)=\sum_{x\in P}G(e,x|r)\rme^{-\sd(e,x)}.
\]

In~\cite[Example~C]{DWY}, we proved that if a finitely generated group $\Gamma_0$ can be endowed with a symmetric finitely supported admissible probability measure $\mu_0$ such that $\Theta_{\Gamma_0}(r_0,\omega_{\Gamma_0}(r_0))$ is convergent for some $r_0<R_0$, then the free product $\Gamma=\Gamma_0*\Z^d$ can also be endowed with a symmetric finitely supported admissible probability measure $\mu$ such that $\Theta_{\Gamma}(r,\omega_{\Gamma}(r))$ is convergent for some $r\leq R$ depending on $r_0$. Here, $R_0$ denotes the inverse of the spectral radius of $\mu_0$ and $R$ the inverse of the spectral radius of $\mu$.

Note that in this situation, $\Gamma$ is relatively hyperbolic with respect to the conjugates of $\Gamma_0$ and $\Z^d$. In the sequel, when it is considered as a maximal parabolic subgroup of $\Gamma$, we will write $P$ for $\Gamma_0$, to be consistent with the notations of the previous sections. The question of whether such a couple $(\Gamma_0,\mu_0)$ exists was left unanswered, but we announced that it was possible to construct one. We provide the details of this construction here and prove in fact a more precise result.

\medskip
Let $F_n$ be the free group with $n$ generators and let $\Gamma_0=F_n\times F_m$, $n\neq m$. The Cayley graph of $\Gamma_0$ is a Cartesian product of two regular trees $T_{l_1}$, $T_{l_2}$ {with respective degrees $l_1 = 2n$ and $l_2 = 2m$}. We consider the measure $\mu_0$ on $\Gamma_0$ defined by~\eqref{defmubitree}. We deduce the following from Theorem~\ref{thmbitree}.

\begin{prop}\label{transitionphaseparabolicsubgroup}
There exists $1<r_0 <R_0$ such that for $r\leq r_0$, the Poincaré series $\Theta_{\Gamma_0}(r,\omega_{\Gamma_0}(r))$ is divergent and for $r_0<r<R_0$, it is convergent.
\end{prop}

We now recall how the measure $\mu$ is constructed on the free product $\Gamma=\Gamma_0*\Gamma_1$ where $\Gamma_1=\Z^d$. Let $\mu_1$ be any finitely supported symmetric admissible probability measure on $\Z^d$. Following~\cite{DWY}, we assume that $d\geq3$ for convenience, so that the random walk associated with $\mu_1$ is transient at the spectral radius, i.e.the Green function $G_{\mu_1}(e,e|R_1)$ is finite, where $R_1$ is the inverse of the spectral radius of $\mu_1$.

For $0\leq \alpha\leq 1$, we set
\[
\mu_\alpha=\alpha \mu_1+(1-\alpha)\mu_0.
\]
For simplicity, we write $\mu=\mu_\alpha$ below. We write $G$ for the Green function associated with $\mu$ and $G_{i}$ for the Green function associated with $\mu_i$. By~\cite[(10), (11)]{DWY}, there exist two numbers $w_{0,\alpha,r}$ and $w_{1,\alpha,r}$ and continuous non-decreasing functions $\zeta_{i,\alpha}$ of $r\leq R$ such that
\begin{align}
G(e,x|r) &=\frac{1}{1-w_{0,\alpha,r}}G_0(e,x|\zeta_{0,\alpha}(r)),
\; x\in \Gamma_0\label{equationG0}\\
\intertext{and}
G(e,x|r)
&=\frac{1}{1-w_{1,\alpha,r}}G_1(e,x|\zeta_{1,\alpha}(r)),
\; x\in \Gamma_1=\Z^d.\label{equationG1}
\end{align}
Furthermore, we have $\zeta_{\alpha, 0}(r) = \frac{(1 - \alpha) r}{1 - w_{0,\alpha,r}}$, and
\begin{equation}\label{wralpha}
w_{0, \alpha,r} = \sum_{n \ge 1} \bfP \left(X_n = e, X_k \neq e, 1 \leq k < n, \text{ first step chosen using } \alpha \mu_1 \right) r^n.
\end{equation}
Similar expressions hold for $w_{1,\alpha,r}$ and $\zeta_{\alpha,1}(r)$.

\begin{lemm}
For fixed $\alpha$, the functions $w_{i,\alpha,r}$ and $\zeta_{\alpha,i}(r)$ are (strictly) increasing in $r$.
\end{lemm}

\begin{proof}
The functions $w_{i,\alpha,r}$ are power series in $r$ with positive coefficients, so they are increasing. It follows from the above expression of $\zeta_{\alpha,i}$ that these functions are increasing too.
\end{proof}



We prove here the following.
\begin{theo}\label{transitionphasePoincare}
If $\alpha$ is small enough, then there exist $r_*(\alpha)<r_\sharp(\alpha) <R$ so that the following holds. For $r\leq r_*(\alpha)$, the Poincaré series $\Theta_{\Gamma}(r,\omega_{\Gamma}(r))$ is divergent and $\Gamma$ has a parabolic gap for the Green function. On the other hand, at $r=r_\sharp(\alpha)$, it is convergent and $\omega_P(r)=\omega_\Gamma(r)$, where $P=\Gamma_0$.
\end{theo}


In the proof of Theorem~\ref{transitionphasePoincare}, we shall also use the following result which is an enhanced version of~\cite[(14)]{DWY}.

\begin{lemm}\label{lemmaDWY(14)}
For every $\epsilon>0$, there exists $\alpha_0$ such that for $\alpha\leq \alpha_0$, the following holds. For every $r\in [0,R]$ and for every $x\in \Gamma_1\setminus\{e\}$, we have
\[
G(e,x|r)\leq \epsilon.
\]
\end{lemm}

\begin{proof}
By~\cite[Lemma~3.15]{DWY}, there exists $\alpha_0$ such that for $\alpha\leq \alpha_0$, $w_{1,\alpha,r}$ stays bounded away from 1 and moreover, $\zeta_{1,\alpha}(r)$ converges to 0 as $\alpha$ tends to 0 uniformly over $r\leq R$. We conclude from~\eqref{equationG1} that for small enough $\alpha$, independently of $r$, for every $x\neq e\in \Gamma_1$ we have
\begin{equation*}\label{thirdboundpoincare}
G(e,x|r)\leq \frac{G_1(e,x|\zeta_{1,\alpha}(r))}{1-w_{1,\alpha,r}}\leq \epsilon. \qedhere
\end{equation*}
\end{proof}



We are ready to complete the proof of Theorem~\ref{transitionphasePoincare}. Let us first explain briefly how the quantities $\alpha$, $r_*(\alpha)$ and $r_\sharp(\alpha)$ are chosen.

By~\cite[Lemma~3.14]{DWY}, as $\alpha$ converges to 0, $\zeta_{0,\alpha}(R)$ converges to $R_0$. Let us now fix any $r_1>r_0$, where $r_0$ is given by Proposition~\ref{transitionphaseparabolicsubgroup}. Then, there exists $\alpha_1$ such that $\zeta_{\alpha, 0}(R)>r_1$ for any $\alpha\le \alpha_1$. Since $\zeta_{0,\alpha}(r)$ is increasing in $r$, there exist $r_*<r_\sharp <R$ depending on $\alpha$ such that $\zeta_{\alpha,0}(r_*)=r_0$ and $\zeta_{0,\alpha}(r_\sharp)=r_1$.

We will also need to choose some $\epsilon>0$ only depending on $\mu_0$, $\mu_1$ and $r_1$ such that Equation~\eqref{equationepsilonpoincare} holds below. Then, we choose $\alpha$ small enough such that the conclusions of Lemma~\ref{lemmaDWY(14)} holds for such $\epsilon$ and such that there exist $r_*<r_\sharp <R$ with $\zeta_{\alpha,0}(r_*)=r_0$ and $\zeta_{0,\alpha}(r_\sharp)=r_1$.

\begin{proof}[{Proof of Theorem~\ref{transitionphasePoincare}}]
The proof follows the lines of~\cite[Example~C]{DWY}. We~will write as announced above $P$ for $\Gamma_0$ when it is considered as a maximal parabolic subgroup of $\Gamma$. In particular, we write
\[
\omega_{\Gamma_0}(s)=\limsup_{n\to\infty} \frac{1}{n}\log \sum_{x\in \Gamma_0,|x|=n}G_0(e,x|s)
\]
for the growth rate of the Green function $G_0$ associated with $\mu_0$ and
\[
\omega_P(r)=\limsup_{n\to\infty} \frac{1}{n}\log \sum_{x\in P,|x|=n}G(e,x|r)
\]
for the growth rate of the Green function $G$ associated with $\mu$, induced on $P$. By~\eqref{equationG0}, we see that
\begin{equation}\label{equationgrowthrates}
\omega_P(r)=\omega_{\Gamma_0}(\zeta_{0,\alpha}(r)).
\end{equation}
We will also write $\Theta_{\Gamma_0}(r,s)$ for the Poincaré series associated with $\mu_0$ on $\Gamma_0$ and $\Theta_P(r,s)$ for the Poincaré series associated with $\mu$ induced on $P$, defined as above by
\[
\Theta_P(r,s)=\sum_{x\in P}G(e,x|r)\rme^{-s|x|}.
\]

If $\alpha$ is small enough, there exist $r_*<r_\sharp <R$ with $\zeta_{0,\alpha}(r_*)=r_0$ and $\zeta_{0,\alpha}(r_\sharp)=r_1$, where $r_1$ is fixed in $(r_0,R_0)$. First, if $r\leq r_*$, then $\zeta_{0,\alpha}(r)\leq r_0$. By Proposition~\ref{transitionphaseparabolicsubgroup}, $\Theta_{\Gamma_0}(\zeta_{0,\alpha}(r),\omega_{\Gamma_0}(\zeta_{0,\alpha}(r)))$ diverges, hence $\Theta_P(r,\omega_{P}(r))$ also diverges. Consequently, according to~\cite[Corollary~3.9]{DWY}, $\omega_P(r)<\omega_\Gamma(r)$. Moreover, since $\Gamma_1$ is amenable, we also deduce from Proposition~\ref{propamenableparabolic} that the induced growth rate on $\Gamma_1$ considered as a maximal parabolic subgroup is smaller than $\omega_\Gamma(r)$. In other words, $\Gamma$ has a parabolic gap for the Green function on $(0,r_*]$. Therefore, \cite[Lemma~3.7]{DWY} yields that the Poincaré series $\Theta_{\Gamma}(r,\omega_\Gamma(r))$ is divergent for $r\leq r_*$.

Next, fix $r$ at the value $r=r_\sharp$ so that $\zeta_{0,\alpha}(r)=r_1>r_0$. By~\cite[(15)]{DWY},
\[
\Theta_{\Gamma}(r,s)\leq G(e,e|r)\sum_{k\ge 0}\left (\sum_{x\in \Gamma_0\setminus\{e\}}\frac{G(e,x|r)}{G(e,e|r)}\rme^{-s|x|}\right)^k\left (\sum_{y\in \Gamma_1\setminus\{e\}}\frac{G(e,y|r)}{G(e,e|r)}\rme^{-s|y|}\right)^k
\]
By~\eqref{equationgrowthrates}, $\omega_P(r)=\omega_{\Gamma_0}(\zeta_{0,\alpha}(r))=\omega_{\Gamma_0}(r_1)$ Moreover, by Theorem~\ref{thmbitree}, $r_0>1$, so that $r_1>1$ and so we deduce from Lemma~\ref{lowerboundgrowthrateat1} and Proposition~\ref{omegaGcont} that $\omega_{\Gamma_0}(r_1)\geq c>0$. Therefore,
\begin{equation}\label{equationgrowthratepositive}
\omega_{\Gamma}(r)\geq \omega_P(r) = \omega_{\Gamma}(r)\geq \omega_{\Gamma_0}(\zeta_{0,\alpha}(r)) \geq c>0.
\end{equation}
Using~\eqref{equationG0} and Proposition~\ref{transitionphaseparabolicsubgroup}, we see that
\begin{equation}\label{firstboundpoincare}
\sum_{x\in \Gamma_0\setminus\{e\}}\frac{G(e,x|r)}{G(e,e|r)}\rme^{-\omega_\Gamma(r)|x|}\leq \frac{1}{G_0(e,e|r_1)}\Theta_{\Gamma_0}(r_1,\omega_{\Gamma_0}(r_1))
<\infty.
\end{equation}
Also, $\Z^d$ has polynomial growth and so by~\eqref{equationgrowthratepositive}
\begin{equation}\label{secondboundpoincare}
\sum_{y\in \Gamma_1\setminus\{e\}}\frac{1}{G(e,e|r)}\rme^{-\omega_\Gamma(r)|y|}\leq \frac{1}{G(e,e|r)}\sum_{y\in \Gamma_1\setminus\{e\}}\rme^{-c|y|}\leq C_1.
\end{equation}
We then choose $\epsilon>0$ such that
\begin{equation}\label{equationepsilonpoincare}
C_2:=\epsilon\frac{1}{G_0(e,e|r_1)}\Theta_{\Gamma_0}(r_1,\omega_{\Gamma_0}(r_1))C_1<1.
\end{equation}
Note that $\epsilon$ does not depend on $\mu$ but only on $C_1$, $\mu_0$ and $r_1$. In particular, it does not depend on $\alpha$ and so by Lemma~\ref{lemmaDWY(14)}, we can choose $\alpha$ small enough so that for every $y\in \Gamma_1\setminus \{e\}$,
\[
G(e,y|r)\leq \epsilon.
\]


Combining~\eqref{firstboundpoincare}, \eqref{secondboundpoincare}, \eqref{equationepsilonpoincare} and~\eqref{thirdboundpoincare}, we get
\begin{align*}
\Theta_{\Gamma}(r,\omega_\Gamma(r))&\leq G(e,e|r)\sum_{k\ge 0}\left(\epsilon\frac{1}{G_0(e,e|\zeta_{0,\alpha}(r))}\Theta_{\Gamma_0}\bigl(\zeta_{0,\alpha},\omega_{\Gamma_0}(\zeta_{0,\alpha}(r))\bigr)C_1\right)^k\\
&\leq G(e,e|r)\sum_{k\ge 0}C_2^k,
\end{align*}
so that $\Theta_{\Gamma}(r,\omega_\Gamma(r))$ is finite.

Finally, we deduce from~\cite[Lemma~3.7]{DWY} that at $r=r_\sharp$, $\Gamma$ does not have a parabolic gap for the Green function. Since $\Gamma_1$ is amenable, we necessarily have $\omega_\Gamma(r)=\omega_P(r)$ by Proposition~\ref{propamenableparabolic}. This concludes the proof of Theorem~\ref{transitionphasePoincare}.
\end{proof}

\begin{rema}\label{remarkchoiceofr1}
In the proof of Theorem~\ref{transitionphasePoincare}, note that we can choose $r_1$ arbitrarily close to $r_0$. Unfortunately, as $r_1$ goes to $r_0$, $\Theta_{\Gamma_0}(r_1,\omega_{\Gamma_0}(r_1))$ tends to infinity, so $\epsilon$ satisfying~\eqref{equationepsilonpoincare} converges to 0. Consequently, the parameter $\alpha$ also tends to 0. In other words, as $r_1$ tends to $r_0$, we need to choose a measure $\mu$ that tends to the measure $\mu_0$ distributed on $\Gamma_0$. Now, since the functions $\zeta_{i,\alpha}$ of $r$ depend on $\alpha$, we cannot guarantee that $r_{\sharp}$ tends to $r_*$. In particular, we cannot prove that there is a true phase transition for the convergence of the Poincaré series $\Theta_\Gamma(r,\omega_\Gamma(r))$ at $r_*$. However, since $r_1$ can be chosen arbitrarily close to $r_0$ there is a phase transition at~$r_0$ at the level of the parabolic subgroup.
\end{rema}

\begin{coro}\label{corosubmultiplicativityfails}
If $\alpha$ is small enough, then the following holds. For $r=r_\sharp(\alpha)$, $\Gamma$ does not have a parabolic gap for the Green function and~\eqref{roughmultiplicativity} and~\eqref{purelyexponentialgrowth} do not hold.
\end{coro}


\begin{proof}
By Theorem~\ref{transitionphasePoincare}, if $r=r_\sharp(\alpha)$, then the Poincaré series $\Theta_\Gamma(r,\omega_\Gamma(r))$ is convergent and $\Gamma$ does not have a parabolic gap. Assume by contradiction that there exists $C$ such that
\[
H_r(m+n)\leq C H_r(n)H_r(m).
\]
Then, the quantity $CH_r(n)$ is sub-multiplicative, hence by Fekete's lemma,
\[
\omega_\Gamma(r)=\lim_{n\to\infty} \frac{1}{n}\log CH_r(n)=\inf_{n\ge 1} \frac{1}{n}\log CH_r(n).
\]
Thus, for every $n$, we have
\[
CH_r(n)\geq \rme^{n\omega_\Gamma(r)}.
\]
This implies that $\Theta_\Gamma(r,\omega_\Gamma(r))$ diverges, which is a contradiction. In particular, we see that~\eqref{roughmultiplicativity} fails. Now, \eqref{roughmultiplicativity} is a direct consequence of~\eqref{purelyexponentialgrowth}, with $C=(C')^3$, hence~\eqref{purelyexponentialgrowth} also fails.
\end{proof}

We also deduce the following from Theorem~\ref{divergencetypeandatoms} and Theorem~\ref{transitionphasePoincare}.
\begin{coro}\label{transitionphasePattersonSullivan}
If $\alpha$ is small enough, then the following holds. For $r\leq r_*(\alpha)$, the measure $\nu_e$ on the Bowditch boundary has no atom and is supported on the set of conical limit points. For $r=r_\sharp(\alpha)$, it is purely atomic and is supported on the set of parabolic limit points.
\end{coro}


\begin{rema}\label{remarkCGM}
In~\cite{CandelleroGilchMuller}, the authors prove that in the context of free products, we always have $\omega_P(r)<\omega_\Gamma(r)$ for every maximal parabolic subgroup $P$, i.e. $\Gamma$~has a parabolic gap for the Green function. However, their proof relies on an unproved statement, namely that the quantity $H_r(n)$ is sub-multiplicative for every finitely generated group and then apply this property to the maximal parabolic subgroup~$P$, see precisely the proof of~\cite[Lemma~4.7]{CandelleroGilchMuller} and also~\cite[Remark~3.17]{DWY}. However, by Theorem~\ref{thmbitree}, we see that sub-multiplicativity fails for the Cartesian product of two regular trees if $r\geq r_0$. Moreover, by Corollary~\ref{corosubmultiplicativityfails}, in the above example, if $r=r_\sharp(\alpha)$, then $\Gamma$ does not have a parabolic gap for the Green function.
\end{rema}



\subsection{The growth tightness property}\label{s:growthtightness}

Let $d$ be a proper left invariant distance on $\Gamma$. The \emph{growth rate} of $\Gamma$ for $d$ is defined as follows:
\[
\delta(\Gamma,d) :=\limsup_{n\to\infty}\frac{\log \sharp\{x\in \Gamma: d(e,x)\le n\}}{n}
\]
A \emph{nontrivial quotient} $\bar \Gamma$ of $\Gamma$ means that the kernel of the canonical projection $\Gamma\to \bar \Gamma$ is an infinite normal subgroup of $\Gamma$. We say that $\Gamma$ is \emph{growth tight} for the distance $d$ if for every nontrivial quotient $\bar \Gamma$ of $\Gamma$, endowed with the quotient distance $\bar d$ from $d$, we have $\delta(\bar \Gamma,\bar d)<\delta(\Gamma,d)$.

Let us assume that $\Gamma$ is a relatively hyperbolic group. Whenever a maximal parabolic group $P$ has growth rate $\delta(P,d)$ strictly less than $\delta(\Gamma, d)$, we say that $\Gamma$ has a parabolic gap along $P$ fir the distance $d$. When $\Gamma$ has parabolic gap along every maximal parabolic subgroup, we say that $\Gamma$ has the parabolic gap property.


Recall that by Lemma~\ref{quasiisometricdistances}, $\fkd_r(x,y)= {\omega_{\Gamma}(r)} |x^{-1}y|+|x^{-1}y|_r$ is quasi-isometric to the word distance for $r<R$. By Lemma~\ref{NewPoincareSeries}, we have $\delta(\Gamma,\fkd_r)=1$. The following result relates the gap property for Green functions to the gap property for the distance $\delta(\Gamma,\fkd_r)$.

\begin{prop}\label{PropTwoPGs}
Let $\Gamma$ be a group endowed with a probability measure $\mu$ such that the $\mu$-random walk is transient at the spectral radius, i.e.$G(e,e|R)$ is finite. Let $A\subset \Gamma$ be any subset. If $\omega_A(r)<\omega_\Gamma(r)$ for some $1<r\le R $, then $\delta(A,\fkd_r)<1$.
\end{prop}

We need the following lemma, which generalizes~\cite[Lemma~3.1]{Tanaka17} in hyperbolic groups.
\begin{lemm}\label{ContinuitythetaLem}
Under the assumption of Proposition~\ref{PropTwoPGs}, for any $\theta\in\R$, $r\leq R$ and $A\subset \Gamma$, define
\[
\omega_{A, r}(\theta) := \limsup_{n \to \infty} \frac{1}{n} \log \sum_{x \in A, |x| = n} \left[G(e, x | r) \right]^\theta.
\]
Then $\omega_{A, r}(\theta)$ is a convex function on $\R$. If $\Gamma$ is a relatively hyperbolic group, then $\omega_{\Gamma, r}(\theta)$ is a true limit.
\end{lemm}
\begin{proof}
Denote $H^\theta_r(n):=\sum_{x \in A, |x| = n} \left[G(e, x | r) \right]^\theta$. For $\theta_0$, $\theta_1 \in \R$ and $0 < t < 1$, by the H\"older inequality,
\[
H_r^{t \theta_0 + (1 - t) \theta_1}(n)
\leq \left(H_r^{\theta_0}(n)\right)^t \left(H_r^{\theta_1}(n)\right)^{1 - t}.
\]
Thus $\omega_{A, r}$ is convex:
\[
\omega_{A, r}(t \theta_0 + (1 - t) \theta_1) \leq t \omega_{A, r}(\theta_0) + (1 - t) \omega_{A, r}(\theta_1).
\]
If $\Gamma$ is a relatively hyperbolic group, the same proof as in~\cite[Lemma~3.2]{DWY} shows (there with $\theta=1$) that for $A=\Gamma$, the sequence $H_r^\theta(n)$ is sub-multiplicative, that is, $H_r^\theta(n+m)\le C H_r^\theta(n) H_r^\theta(m)$ for some $C>0$. Thus, the limit exists by Fekete's lemma.
\end{proof}

As a convex function on $\R$, $\omega_{\Gamma, r}(\theta)$ is a continuous function of $\theta\in \R$, and is differentiable, except maybe at countably infinitely many points.

\begin{proof}[Proof of Proposition~\ref{PropTwoPGs}]
To show $\delta(A,\fkd_r)<1$, it suffices to find some $\epsilon > 0$ such that
\[
\sum_{x \in A} \rme^{- (1 - \epsilon) \fkd_r(e, x)} = \sum_{x \in A} \rme^{- (1 - \epsilon) \omega_{\Gamma}(r) |x|} \left[G(e, x | r) \right]^{1 - \epsilon} < \infty.
\]



By Lemma~\ref{ContinuitythetaLem}, the function $\omega_{A,r}(\theta)$ is continuous in $\theta\in\R$. If $\omega_A(r) < \omega_{\Gamma}(r)$, we can choose $\epsilon,\eta>0$ small enough so that
\[
\omega_A(r) + \eta < (1 - \epsilon) \omega_{\Gamma}(r),
\]
and at the same time, by continuity of $\omega_{A,r}(\theta)$ at $\theta=1$, the following holds: for large enough $n$,
\[
\sum_{x \in A, |x| = n} \left[G(e, x | r) \right]^{1 - \varepsilon} \leq \rme^{n \left(\omega_{A,r}(1) + \eta \right)}.
\]
By definition, $\omega_{A,r}(1)=\omega_A(r)$, so the two inequalities above yield
\[
\sum_{x \in A} \rme^{- (1 - \epsilon) \omega_{\Gamma}(r) |x|} G(e, x | r)^{1-\epsilon} < \infty,
\]
which is the desired inequality.
\end{proof}

Given $f\in \Gamma$ and $\epsilon>0$, let $\clV_{\epsilon, f}$ be the set of \emph{barrier-free} elements $x\in \Gamma$, that is, elements for which the $\epsilon$-neighborhood of some geodesic $[e,x]$ contains a geodesic segment representing $f$. The following is analogous to~\cite[Theorem~C]{YangSCC}.

\begin{lemm}\label{BarrierFreeLem}
Let $\Gamma$ be a relatively hyperbolic group with parabolic gap property for Green function. Then there exists some $\epsilon>0$ such that the set $\clV:=\clV_{\epsilon, f}$ has growth rate strictly less than 1 for any element $f\in \Gamma$: $\delta(\clV, \fkd_r)<1$.
\end{lemm}
\begin{proof}
By Proposition~\ref{PropTwoPGs}, it suffices to prove $\omega_{\clV}(r)<\omega_{\Gamma}(r)$. Set
\[
a^\omega(n) = \rme^{-\omega n}\sum_{x\in \clV, |x|=n} G(e,x;r).
\]
Assume that $\omega_{\clV}(r) >\omega_P(r)$ for every maximal parabolic subgroup $P$; otherwise the parabolic gap concludes the proof. If $\omega_{\clV}(r)>\omega >\omega_P(r)$, one obtains
\[
a^\omega(n+m) \le c_0 \sum_{1\le i\le n} a^\omega(i) \sum_{1\le j\le n} a^\omega(j)
\]
by the same argument of~\cite[Lemma~3.7]{DWY} where $a^\omega(n)$ is summed up over $\Gamma$ instead of $\clV$. This implies via a variant of Fekete's lemma in~\cite[Lemma~4.3]{DPPS} that the series $\sum_{x\in \clV} \rme^{-s |x|}G(e,x;r)$ diverges at $s=\omega_{\clV}(r)$.

Fix any $L>0$. We choose an \emph{$L$-separated net} $A\subset \clV$ for the word distance: for any $x, y\in A$ we have $|x^{-1}y|>L$ and for any $y\in \clV$, there exists $x\in A$ such that $|x^{-1}y|\le L$. Note that if $|x^{-1}y|\le L$, then $G(e, x;r)\asymp_L G(e,y;r)$. Since any ball of radius $L$ in word distance contains a fixed number of elements, we deduce that $\Theta_A(r, s)\asymp_L \Theta_A(r, s)$ whenever they are finite. Thus, $\omega_A(r)=\omega_{\clV}(r)$.

Following~\cite[Section~4.2]{YangSCC}, we use a ping-pong argument to construct a \emph{free product of sets} inside $\Gamma$: if $L$ is large enough, there exist a finite set of elements $B\subset \Gamma$ such that the set $\clW(A, B)$ of alternating words over $A$ and $B$ embeds into $\Gamma$ as a free semi-group under the evaluation map. This construction uses only the word distance. Now, by~\cite[Lemma~3.8]{DWY}, we have $\omega_A(r)<\omega_\Gamma(r)$ and then $\delta(\clV, \fkd_r)<1$ by Proposition~\ref{PropTwoPGs}.
\end{proof}

\begin{prop}\label{growthrightnessandparabolicgap}
If a relatively hyperbolic group $\Gamma$ has parabolic gap for the Green function, then it is growth tight for the distance $\fkd_r$. Otherwise, there exists a nontrivial quotient $\bar \Gamma$ such that $\delta(\bar \Gamma, \bar{\fkd_r})=\delta(\Gamma, \fkd_r)$.
\end{prop}

\goodbreak
\begin{proofc}
\begin{proof}[{(1)}]
We follow the proof of~\cite[Corollary~4.6]{YangSCC} in our setup. Let $N$ be the infinite kernel of $\Gamma \to \bar \Gamma$. We form a set $A$ by choosing a shortest representative $h\in hN$ for each $hN\in \bar G$ so that $\fkd_r(e, h)=\fkd_r(e, hN)$. By definition of the quotient distance, the growth rate of the set $A$ for $\fkd_r$ is exactly the growth rate of $\bar \Gamma$ for $\bar{\fkd_r}$.

We now choose a sufficiently long loxodromic element $f\in N$, which exists since $N$ is infinite. If $|f|/ \epsilon$ is large enough, we see that any geodesic $[e, h]$ cannot contain $f$ in its $\epsilon$-neighborhood. Indeed, if not, the loxodromic element $f$ produces two transition points on some $[e,h]$ with a distance comparable with $|f|$. Now, we use the following fact given by Lemma~\ref{weakAncona}: if $u, v$ are two transition points in this order on a word geodesic $[x,y]$, then
\begin{align}\label{almostEqualityEQ}
\fkd_r(x,y) \simeq \fkd_r(x,u)+\fkd_r(u,v)+\fkd_r(v,y)
\end{align}
where $\simeq$ denote the equality up to a uniform additive constant. We could then shorten $\fkd_r(e,h)$ by an amount $\fkd_r(e,f)=\omega_\Gamma(r)|f|+|f|_r$, giving a contradiction with the above choice of $h\in hN$ as the shortest one.

In other words, we proved that $A\subset \clV_{\epsilon, f}$. Hence, $\delta(A, \fkd_r)\le \delta(\clV_{\epsilon, f},\fkd_r)<1$ by Lemma~\ref{BarrierFreeLem}. The growth tightness follows.
\let\qed\relax
\end{proof}
\begin{proof}[{(2)}]
Assume that $\omega_P(r)=\omega_\Gamma(r)$ for a maximal parabolic subgroup $P$. Then,
\[
\sum_{x \in P}
\rme^{- \fkd_r(e, x)}=\sum_{x \in P} \rme^{-\omega_{P}(r)} G(e,x|r)
\]
hence, we see that the growth rate for $\fkd_r$ induced on $P$ equals $1$.

Fix a loxodromic element $f\in G$. For any large enough $n$, the quotient group $\bar \Gamma$ defined as $G/\llbangle f^n\rrbangle$ is again a relatively hyperbolic group, and $P\cap \llbangle f^n\rrbangle$ is trivial (see~\cite[Lemma~8.9]{YangPS}). Thus, the set of elements in $P$ embeds into $\bar \Gamma$ whose image we denote by $\bar P$, so $\delta(\bar P, \bar{\fkd_r})\ge \delta(P, \fkd_r)=1$. Therefore, $\delta(\bar \Gamma, \bar{\fkd_r})=1$.
\end{proof}\let\qed\relax
\end{proofc}



Relatively hyperbolic groups endowed with a word distance are always growth tight by~\cite{ACT,YangTight}. In fact, any co-compact action of a relatively hyperbolic group on a proper geodesic space contains a contracting element and thus is growth tight. Here, the existence of a contracting element in the co-compact action follows from the fact that in a relatively hyperbolic group, a loxodromic element is contracting with respect to all word quasi-geodesics: any $c$-quasi-geodesic outside the $C$-neighborhood of the axis has $C$-bounded projection for some $C=C(c)$. See~\cite[Proposition~8.5]{GP16}.

On the contrary, as a corollary of Theorem~\ref{transitionphasePoincare} and Proposition~\ref{growthrightnessandparabolicgap}, growth tightness for $\fkd_r$ may fail and depends on $r$.

\begin{theo}\label{transitionphasegrowthtightness}
There exists a relatively hyperbolic group $\Gamma$ endowed with a finitely supported symmetric and admissible probability measure $\mu$ such that the following holds. There exist $1<r_*<r_\sharp<R$ such that $\Gamma$ endowed with the distance~$\fkd_r$ is growth tight for $r\le r_*$, but is not for $r=r_\sharp$.
\end{theo}



Note that the proper distance $\fkd_r$ is quasi-isometric to any word distance for $r<R$ by Lemma~\ref{quasiisometricdistances}. We say that a metric space $(X,d)$ is \emph{$D$-coarsely geodesic} for some\pagebreak{} $D>0$ if for any two points $x,y\in X$, there exists a $(1,D)$-quasi-isometric embedding $\phi: [0, l]\to X$ for $l:=d(x,y)$ so that $\phi(0)=x,\phi(l)=y$, and
\[
\bigl|d(\phi(m), \phi(n))-|m-n|\bigr|\le D
\]
for any $0\le m\le n\le l$. It is an open question whether the Green distance is a geodesic distance on hyperbolic groups, see~\cite[Section~1.7]{BHM11}. We shall however derive the following corollary from Theorem~\ref{transitionphasegrowthtightness}.

\begin{coro}\label{drnotgeodesic}
For $r=r_\sharp$, $(\Gamma,\fkd_r)$ is not a coarsely geodesic metric space.
\end{coro}

The proof requires the following observation of independent interest. Recall that an element of infinite order $g$ in a finitely generated group $\Gamma$ is called contracting for a distance $d$ on $\Gamma$ if any $d$-metric ball in $\Gamma$ disjoint with the subgroup $\langle g\rangle$ has $C$-bounded projection to $\langle g\rangle$ for some universal constant $C>0$. In a $D$-coarsely geodesic metric space, this is equivalent to the bounded image property: there exists $C=C(D)>0$ such that any $D$-coarse geodesic outside the $C$-neighborhood of $\langle g\rangle$ has shortest projection of diameter at most $C$ to it. For simplicity, we can take the same $C$ for both statements.
\begin{lemm}\label{contractingdrmetric}
Any loxodromic element in a relatively hyperbolic group is contracting with respect to $\fkd_r$ where $1\le r\le R$.
\end{lemm}
\begin{proof}
By~\eqref{almostEqualityEQ}, the proper distance $\fkd_r$ is coarsely additive along the set of transition points on the word geodesic. That is, if $z$ is a transition point on $[x,y]$ we have $\fkd_r(x,y)\ge \fkd_r(x,z)+\fkd_r(z,y)-D$ for some universal $D>0$. Let $\gamma$ be a quasi-geodesic preserved by a loxodromic element $g$. Then $\gamma$ is $C_0$-contracting with respect to the word distance for some $C_0>0$. We claim that the shortest projection $z$ of any point $x$ to $\gamma$ for the distance $\fkd_r$ is $D_0$-close to the shortest projection $w$ of $x$ to $\gamma$ for the word distance. Indeed, as $w$ is uniformly close to a transition point on $[x,z]$, we see that $\fkd_r(x,z)+D_0\ge \fkd_r(x,w)+\fkd_r(w,z)$ for some $D_0=D_0(C_0)>0$. By the definition of $\fkd_r$-shortest projection, we have $\fkd_r(x,z)\le \fkd_r(x,w)$ and thus the claim follows.

We now prove that $\gamma$ is contracting for $\fkd_r$. Pick any $\fkd_r$-distance ball $B$ centered at $x$ disjoint with $\gamma$. Let $y\in B$ so that the projections denoted by $u, v$ respectively of $x,y$ to $\gamma$ realize the $\fkd_r$-diameter of the projection of $B$ to $\gamma$. By Lemma~\ref{quasiisometricdistances}, $\fkd_r$ is quasi-isometric to the word distance for $1\le r<R$. For $r=R$, by~\cite[Lemma~2.1]{GouezelLalley}, $f(n):=\max\{G(e,x|R): x\in S_n\}\to 0$ as $n\to \infty$, so $f(n)$ is a proper function. We can thus choose $\fkd_r(u, v)$ large enough so that $|u^{-1}y|\ge C_0$, hence the contracting property of $\gamma$ in word distance implies that $u,v$ are uniformly close to transition points on $[x,y]$. By the additive property of $\fkd_r(x,y)$ along transition points, we obtain $\fkd_r(x,y)\ge \fkd_r(x,u)+\fkd_r(u,v)+\fkd_r(v,y)-D_1$ for some $D_1>0$. As~$B$ is a $\fkd_r$-distance ball disjoint with $\gamma$, we have $\fkd_r(x,y)\le \fkd_r(x,u)$. We then obtain a contradiction if $\fkd_r(u,v)>D_1$. Thus, the contracting property for $\fkd_r$ follows.
\end{proof}

\begin{proof}[{Proof of Corollary~\ref{drnotgeodesic}}]
Assume on the contrary that $\fkd_r$ is a coarsely geodesic distance for $r>r_\star$. By Lemma~\ref{contractingdrmetric}, the axis of any loxodromic element satisfies the bounded image property for coarse geodesics, so any loxodromic element is contracting in the sense of~\cite{YangTight}. As the action on $\Gamma$ is co-compact, the same argument as in~\cite{YangTight} holds verbatim by replacing word geodesics with coarse $\fkd_r$-geodesics and we can show that $\fkd_r$ on $\Gamma$ is growth tight. However, this contradicts Theorem~\ref{transitionphasegrowthtightness}. Thus the corollary follows.
\end{proof}

In~\cite[Question 1]{ACT}, Arzhantseva--Cashen--Tao asked whether growth tightness is invariant among cocompact actions on geodesic metric spaces. Cashen--Tao~\cite{CaTao16} showed the first examples of product groups with growth tightness for one generating set but not for another generating set. Examples of non-growth tight relatively hyperbolic groups with non-cocompact actions were already considered in~\cite[Observation~8.8]{ACT}. As the action is not-cocompact, the induced pseudo-distance on the group pulled back from the action is not quasi-isometric to the word distance. It is natural to ask the following variant of Arzhantseva--Cashen--Tao's question for relatively hyperbolic groups about quasi-isometry invariance of growth tightness: if a proper (pseudo-)distance $d$ on $\Gamma$ is quasi-isometric to the word distance, does the growth tightness for $d$ hold? In our last corollary, we produce examples of non-geodesic distances on relatively hyperbolic groups $\Gamma$ that answer negatively this question.
\begin{coro}\label{corgrowthtightness}
There exists a relatively hyperbolic group with a proper left invariant distance quasi-isometric to the word distance which does not have the growth tightness property.
\end{coro}


\appendix
\section{Probability that the branching random walk visits partial cones}\label{app}


We consider a relatively hyperbolic group $\Gamma$. Our goal is to prove Proposition~\ref{proppartialcone}. The proof basically consists of a reorganization of arguments of~\cite{DWY}. We first recall some notations and definitions.

Given a subset $A$ of $\Gamma$, we write $G(x,y;A|r)$ for the Green function restricted to paths staying in $A$, expect maybe the first and last point. That is,
\[
G(x,y;A|r)=\sum_{n\ge 0}\;\sum_{z_1,\dots,z_{n-1}\in A}r^n\mu\left(x^{-1}z_1\right)\mu\left(z_1^{-1}z_2\right)\dots\mu\left(z_{n-1}^{-1}y\right).
\]

Fix $C>0$ and $x\in \Gamma$. The \emph{$C$-partial cone} $\Omega(x,C)$ consists of points $z\in G$ such that $x$ is within $C$ of an $(\eta, L)$-transition point on the geodesic $[e,z]$. Let $C>0$ be any sufficiently large constant given by~\cite[Lemma~2.9]{DWY} so that the relative thin triangle property holds for $(\eta, L)$-transition points: for every triple $(x,y,z)$ of~$\Gamma$, any $(\eta,L)$-transition point on $[x,y]$ is within $C$ of either an $(\eta,L)$-transition point on $[x,z]$ or an $(\eta,L)$-transition point on $[y,z]$.

Let $B([x,z])$ be the ball centered at the middle point of $[x,z]$ of radius ${d(x,z)/2}$. Define $U(x)$ to be the union of the balls $B([x,z])$ for all geodesics $[x,z]$ between $x$ and $z\in \Omega(x,C)$. That is,
\[
U(x):= \bigcup\bigl\{B([x,z]):\, \forall [x,z], \forall z\in \Omega(x,C)\bigr\}.
\]
It is clear that $\Omega(x, C)$ is contained in $U(x)$.

Fix $\epsilon\in (0, 1/2)$.
Let $U_\epsilon(x)$ be the set of points $z\in U(x)$ such that $[x,z]$ contains a transition point $w$ being at distance at least $\epsilon d(x,z)$ to one of the endpoints:
\[
\max\{d(w, x), d(w,z)\}\ge \epsilon d(x,z).
\]
For any $m\ge 1$, let $U_{\epsilon}(x,m)$ be the set of elements $z\in U_{\epsilon}(x)$ such that $d(x,z)\ge m$.


We now consider a finitely supported symmetric and admissible probability measure $\mu$ on $\Gamma$ and a probability measure $\nu$ on $\bbN$. We denote by $\BRW(\Gamma,\nu,\mu)$ the branching random walk associated with $\nu$ and $\mu$. In what follows, we shall often use the following estimates proved in~\eqref{lowerboundquasiisometry}. There exists $\alpha>0$ such that for any $x\in \Gamma$, for every $r\geq 1$,
\begin{align}\label{GrLBD}
G(e, x|r)\ge \rme^{-\alpha |x|}.
\end{align}
Also, by Lemma~\ref{purelyexponentialgrowth} there exists $c_1>0$ such that for any $n\ge 1$,
\begin{align}\label{SnGrowth}
\sharp S_n\le c_1 e^{vn}
\end{align}



\begin{lemm}\label{TraceAvoidTransBall}
For any $\epsilon\in (0,1/2)$, there exists $\kappa_0>0$ such that for every $\kappa\ge \kappa_0$, the following holds. For all but finitely many $x\in \Gamma$: the following event
\begin{align*}
E_1:=\{\BRW(\Gamma,\nu,\mu)~\text{first enters}~U(x)~\text{at a point}~z\in U_{\epsilon}(x,\kappa \log |x|)\}
\end{align*}
has probability at most $G(e,x|r)$.
\end{lemm}

\begin{proof}
Let us freeze all particles of $\BRW(\Gamma,\nu,\mu)$ when the event $E_1$ happens, and denote by $\clZ$ the collection of frozen particles. Set $m:=\kappa \log n$ for simplicity, where $n=|x|$. Then for $z \in \clZ$ we have
\begin{enumerate}
\item $d(x,z)\ge m$,
\item $\max\{d(y,x), d(y,z)\}>\epsilon d(x,z)-3C$ where $y$ is an $(\eta,L)$-transition point on $[e,z]$ given by~\cite[Lemma~6.6]{DWY}.
\end{enumerate}
As the genealogy path from $e$ to $z$ does not intersect $B(y,\epsilon d(x,z)-3C)$, the expected number of particles frozen at $z\in U_\epsilon(x)$ is upper bounded by
\[
G(e,z; [U_\epsilon(x)]^c|r) \le G\bigl(e,z; [B(y,\epsilon d(x,z)-3C)]^c|r\bigr) \le e^{- e^{\delta [\epsilon d(x,z)-3C]}}.
\]
where $\delta=\delta(\eta, L)$ be given by~\cite[Proposition~3.5]{DG21}.
As a consequence, there exist $\epsilon_1=\epsilon_1(\epsilon,\delta,v)$ and $n_0>0$ such that for any $m>\kappa\log n_0$, we have by~\eqref{SnGrowth} that
\[
\bfE[\sharp \clZ]\le \sum_{z\in U_\epsilon(x,m)} G\bigl (e,z; [B(y,\epsilon d(x,z)-3C]^c|r\bigr) \le \sum_{k=m}^{\infty} c_1\cdot \rme^{kv- e^{\delta(\epsilon k-3C)}} \le e^{-e^{\epsilon_1 m}}.
\]
Choose $\kappa$ so that $n^{\kappa \epsilon_1} >\alpha n$ holds for any $n>n_0$. Then,
\[
\bfP(E_1) \le \bfE(\sharp \clZ\ge 1) \le e^{-n^{\epsilon_1\kappa}} \le \rme^{-\alpha n} \le G(e,x|r)
\]
where the last inequality uses~\eqref{GrLBD}.
\end{proof}



Similarly, we prove the following.
\begin{lemm}\label{TraceAvoidTransBall2}
For every $K\geq 1$ and $\hat C\geq 0$, there exists $\kappa_0>0$ such that for all $\kappa\geq \kappa_0$, the following holds. For any sufficiently large $n\ge 1$, the following event
\pagebreak
\[
E_2:=\left\{
\begin{aligned}
&\BRW(\Gamma,\nu,\mu)~\text{eventually visits a point}~z\in \Omega(x,\hat C) \text{ with}\\
& d(e,z)\leq Kd(e,x)\text{ but without entering }~B(x,\kappa \log |x|) \text{ where } |x|\ge n
\end{aligned}
\right\}
\]
has probability at most $G(e,x|r).$
\end{lemm}

\begin{proof}
We freeze particles when the event $E_2$ happens and denote by $\clZ$ the set of frozen particles. By~\cite[Lemma~2.11]{DWY}, if $y\in [e,z]$ is a transition point $\hat C$-close to $x$, the expected number of particles frozen at $z\in B(e,K|x|)$ is upper bounded by
\[
G\left(e,z; \left[B(y,\kappa \log |x|-\hat C)\right]^c|r\right) \le e^{- n^{\delta \kappa }}
\]
where $c_2$ depends on $\hat C, \delta$. Thus, we have
\[
\bfE[\sharp \clZ]\le \sum_{m\ge n}\;\sum_{z\in S_{Km}} G\left(e,z; \left[B(y,\kappa \log m-\hat C)\right]^c|r\right) \le \sum_{m\ge n} c_1 e^{vKm} e^{- c_2m^{\delta\kappa}}.
\]
Choose $\kappa, n_0>0$ so that $c_2 m^{\delta\kappa}-vKm>\alpha m$ holds for any $m\ge n_0$. The conclusion follows again.
\end{proof}

\begin{lemm}\label{exponentialmoments}
For any $K>1$, there exist $\epsilon_0, \kappa_0$ such that for all $\epsilon\leq \epsilon_0$ and $\kappa\geq \kappa_0$, the following holds. There exists $c<0$ such that for all but finitely many $x\in \Gamma$, the following event
\[
E_3:=\left\{
\begin{multlined}
\BRW(\Gamma,\nu,\mu)~\text{first enters}~U(x)~\text{at a point}\\
z\in U(x)\setminus U_{\epsilon}(x,\kappa \log |x|)~\text{with}~|z|\ge K|x|
\end{multlined}
\right\}
\]
has probability at most $c\cdot G(e,x|r)$.
\end{lemm}



\begin{proof}
Let $V$ be the set of $z\in U(x)\setminus U_{\epsilon}(x,\kappa \log |x|)$ satisfying $|z|\ge K|x|$. By definition,
\begin{enumerate}
\item\label{prflemmA.3.1} either $d(z, x)\le \kappa \log |x|$,
\item\label{prflemmA.3.2} or the $[\epsilon, 1-\epsilon]$-percentage of $[x,z]$ does not contain any $(\eta,L)$-transition point.
\end{enumerate}
If $K $ and $ \kappa$ are fixed, noticing that
\[
(K-1)|x|\le d(x,z)\le (1+1/K)|z|,
\]
the case~\eqref{prflemmA.3.1} is impossible for sufficiently large $|x|$. Thus, it suffices to consider the case~\eqref{prflemmA.3.2}. Set $K_1=\epsilon(1+1/K)$ and $K_2=(1-2\epsilon)(1-1/K)$. By~\cite[Lemma~6.1]{DWY}, there exist a unique coset $P_z\in \bbP$ such that if $y_1,y_2$ are the entrance and exit points of $[x,z]$ in $N_\eta(P_z)$, then
\begin{align}
\max\{d(x,y_1), d(y_2,z)\} &\le \epsilon d(x,z)\le \epsilon(1+1/K)|z|\le K_1|z|,\label{y1y2Def}
\\
\intertext{and}
d(y_1,y_2)\geq (1-2\epsilon)d(x,z) &\ge (1-2\epsilon)(1-1/K)|z|\ge K_2|z|.\label{y1y2EQ}
\end{align}


Before moving on, we need the following facts about $y_1$. By relative thin triangle property, there exists a constant $\hat C_1$ depending only on $C$ so that $x$ is $\hat C$-close to a~transition point on $[e,y_1]$. Moreover, there exists $\hat C_2$ depending on $C, \eta$ so that the projection $\pi_{N_{\eta}(P_z)}(e)$ of $e$ to $N_{\eta}(P_z)$ is within $\hat C_2$ of $y_1$ and $y_1$ is $\hat C_2$-close to a~transition point on $[e,w]$ for any $w\in N_\eta(P_z)$. For given $P\in \bbP$, let $P(y_1)$ denote the set of $w\in N_\eta(P)$ with $d(w,y_1)\leq d(e,y_1)$, where $y_1$ is $\hat C_2$-close to $\pi_{N_\eta(P)}(e)$.

Consider first the sub-event $E_{30}$ of $E_3$, where $\BRW(\Gamma,\nu,\mu)$ enters $N_\eta(P_z)$ at a point $w\in P_z(y_1)$. Thus, $d(e,w)\le d(w,y_1)+ d(e,y_1)\le 2|y_1|$.
The same proof of~\cite[Lemma~6.6]{DWY} implies that $B(y_1, d(y_1,x)-3C)$ is contained in $U(x)$. Thus, the particle does not visit $B(y_1, d(y_1,x)-3C)$. Assume first that $|y_1|>2|x|$. Note that there exists $n_0, \kappa_0>0$ so that for any $|x|\ge n_0$, we have
\[
\kappa_0 \log |y_1| \le |y_1|-|x|-3C\le d(y_1,x)-3C.
\]
In particular, the branching random walk does not visit $B(y_1,\kappa_0 \log |y_1|)$ before arriving at $w$. Now, if $|y_1|\le 2|x|$, then $d(e,w)\le 2|y_1|\le 4|x|$. By definition of $E_3$, the branching random walk does not visit $B(x,\kappa_0 \log |x|)$. In summary, this sub-event $E_{30}$ is included into the event $E_2$ in Lemma~\ref{TraceAvoidTransBall2} with constants $\hat C$ and $K=4$, so there exists $\kappa_0\ge \kappa>0$ so that the probability of $E_{30}$ is at most $G(e,x|r)$.

Now, it remains to consider the particles of $\BRW(\Gamma,\nu,\mu)$ in the event $E_3$ that do not enter $N_\eta(P_z)$ at some point $w\in P_z(y_1)$. Then, we have the following two sub-events denoted by $E_{31}$ and $E_{32}$ respectively: the particles either do not visit $N_\eta(P_z)$ at all or do visit $N_\eta(P_z)$ but at a first entrance point $w$ not in $P_z(y_1)$. Let us denote by $W$ the set of points $w\in N_\eta(P)\setminus P(y_1)$ for all $P\in \bbP$ where $y_1$ is $\hat C_2$-close to $ \pi_{N_\eta(P)}(e)$ and $|y_1|>|x|$. In the first case, we freeze particles when they first enter the set $V$. In the second case, we freeze particles when they first enter the set $W$. We denote by $\clZ_{31}$ and $\clZ_{32}$ respectively the sets of frozen particles. We have
\begin{align}
\bfE[\sharp \clZ_{31}]&\leq \sum_{z\in V} G\bigl(e,z;[N_\eta(P_z)]^c|r\bigr),\label{Z31}\\
\bfE[\sharp \clZ_{32}]&\leq \sum_{w\in W} G\bigl(e,w;[N_\eta(P_z)]^c|r\bigr),\label{Z32}
\end{align}


We first bound $\bfE[\sharp \clZ_{32}]$ in~\eqref{Z32}. Recall that $\pi_{N_{\eta}(P_z)}(e)$ is $\hat C_2$-close to $y_1$. If the branching random walk $\BRW(\Gamma,\nu,\mu)$ enters $N_\eta(P_z)\setminus P_z(y_1)$ at a point $w$, by~\cite[Lemma~2.12]{DWY}, for every $M\geq 0$, there exists $\eta_0$ such that for all $\eta\geq \eta_0$,
\[
G\bigl(e,w;[N_\eta(P_z)]^c|r\bigr)\leq \rme^{-M d\left(\pi_{N_{\eta}(P_z)}(e),w\right)}\leq c_0 \rme^{-M d(w,y_1)}
\]
for some $c_0=c_0(\hat C_2)>0$. Using~\eqref{GrLBD}, we first sum up over $y_1$ with $|y_1|>|x|$ and then $w\in P_z(y_1)$ with $d(y_1, w)>d(e, y_1)$:
\[
\bfE[\sharp \clZ_{32}]\leq \sum_{n\ge |x|} c_0 \cdot \rme^{vn} \sum_{m\ge n} \rme^{(v- M)m}.
\]
Choosing $M>2v+\alpha$ that is independent of $\eta$, we have that for $\eta\geq \eta_0$,
\[
\bfE[\sharp \clZ_{32}]\leq c\rme^{-\alpha |x|} \le c G(e,x|r)
\]
where $c$ depends on $c_0$, thus on $\eta$.

We are left to bound $\bfE[\sharp \clZ_{31}]$ in~\eqref{Z31}. As the support of $\mu$ is finite, we can replace each edge in the geodesic from $z$ to $y_2\in N_\eta(P_z)$ by a $\mu$-trajectory with uniformly bounded length. Moving possibly the endpoint $y_2$ up to a bounded distance depending on $\supp(\mu)$, this produces a trajectory outside $N_\eta(P_z)$ for the $\mu$-random walk from $z$~to $y_2$ so that its length is bounded above by a linear function of $d(y_2,z)$. This implies the existence of a positive $\beta$ independent on $z, y_1,y_2$ and $\eta$ such that
\[
G\bigl(z,y_2;[N_\eta(P_z)]^c|r\bigr)\geq \rme^{-\beta d(y_2,z)}\ge \rme^{-\beta K_1 |z|}.
\]
Taking into account that
\[
G\bigl(e,z;[N_\eta(P_z)]^c|r\bigr) \cdot G\bigl(z,y_2;[N_\eta(P_z)]^c|r\bigr) \le G(e,e|r) \cdot G\bigl(e,y_2;[N_\eta(P_z)]^c|r\bigr)
\]
we obtain
\begin{equation}\label{proofexponentialmoments}
G\bigl(e,z;[N_\eta(P_z)]^c|r\bigr)\leq G\bigl(e,y_2;[N_\eta(P_z)]^c|r\bigr)\cdot \rme^{\beta K_1 |z|},
\end{equation}
As above, the projection $\pi_{N_{\eta}(P_z)}(e)$ has a distance at most $\hat C_2$ depending on $\eta$ to $y_1$. By~\cite[Lemma~2.12]{DWY}, for every $M\geq 0$, there exists $\eta_0$ such that for all $\eta\geq \eta_0$,
\[
G\bigl(e,y_2;[N_\eta(P_z)]^c|r\bigr)\leq \rme^{-M d\left(\pi_{N_{\eta}(P_z)}(e),y_2\right)}\leq c_0 \rme^{-M d(y_2,y_1)}\le c_1\rme^{-K_2 M|z|}
\]
for some $c_0=c_0(\eta)>0$.
Summing over $z\in V$ with $|z|>K|x|$, choosing $M>0$ so that $M K_2>\beta K_1+v+\alpha$, we have by~\eqref{proofexponentialmoments} and~\eqref{GrLBD} that for $\eta\geq \eta_0$,
\[
\bfE[\sharp \clZ_{31}]\leq \sum_{n\ge K|x|} c_0c_1 \cdot \rme^{(v+\beta K_1- M K_2)n}\le c\rme^{-\alpha |x|} \le c G(e,x|r).
\]
where $c$ depends on $c_0, c_1$. The Lemma~\ref{exponentialmoments} is proved.
\end{proof}

We can now finish the proof of Proposition~\ref{proppartialcone}.
\begin{proof}[Proof of Proposition~\ref{proppartialcone}]
Using the thin-triangle property, we see that if $z$ is contained in $U(x)\setminus U_{\epsilon}(x,\kappa \log |x|)$, then $z\in \Omega(x, \hat C)$ for some uniform $\hat C$.

Fix any $K>1$. Let $\epsilon$ be small enough and $\kappa$ be large enough such that the conclusions of Lemmas~\ref{TraceAvoidTransBall}, \ref{TraceAvoidTransBall2} and~\ref{exponentialmoments} hold. Using~\eqref{GrLBD}, the probability that the branching random walk visits $B(x, \kappa \log |x|)$ is bounded by
\[
\sum_{z \in B(x, \kappa \log |x|)} G(e, z|r)\le c\rme^{v\kappa \log |x|} G(e,x|r) \rme^{\alpha \kappa \log |x|}\le |x|^{\beta}G_r(e,x)
\]
for some $\beta$ depending on $\alpha, \kappa$ and $v$.

Now assume that the branching random walk visits $U(x)$ through paths outside $B(x, \kappa \log |x|)$. If the events $E_1$ and $E_3$ do not happen, then the event $E_2$ happens: $\BRW(\Gamma,\nu,\mu)$ first visits $U(x)$ {at a point} $z\in U(x)\setminus U_{\epsilon}(x,\kappa \log |x|)$ with $|z|< K|x|$, without entering $B(x, \kappa \log |x|)$. Therefore,
\begin{multline*}
\bfP\bigl(\BRW(\Gamma,\nu,\mu)\text{ visits } \Omega(x, C)\bigr)\\
\begin{aligned}
&\leq \bfP \bigl(\BRW(\Gamma,\nu,\mu)\text{ visits } U(x)\bigr) \\
&\leq \bfP\bigl(\BRW(\Gamma,\nu,\mu)\text{ visits }B(x, \kappa \log |x|)\bigr)+\bfP(E_1\cup E_2\cup E_3) \\
&\leq C\left(1+|x|^{\beta}\right)G(e,x|r) \qedhere
\end{aligned}
\end{multline*}
\end{proof}

\let\OLDthebibliography\thebibliography
\renewcommand\thebibliography[1]{
\OLDthebibliography{#1}
\setlength{\parskip}{0pt}
\setlength{\itemsep}{7pt plus 0.3ex}}

\bibliography{dussaule}
\end{document}
